\part{Root algebras, extension geometry and operators}
\chapter{Lattices and automorphic denominators}
\label{comp:ch:root-data}

The denominator records signed dimensions. A bracket requires the
labelled generators and relations behind those dimensions. The
lattice fixes their degrees, the automorphic product fixes the
chamber correction, and that correction specifies the minimal
Borcherds presentation.

\section{The Lorentzian root lattice}
Let
\[
 L=\Z f_2\oplus\Z f_3\oplus\Z f_{-2}
\]
with bilinear form
\[
 (f_2,f_{-2})=-1,
 \qquad (f_3,f_3)=2,
\]
and all remaining basis pairings zero. Its signature is $(2,1)$. Define the type-II root lattice
\[
 L_{II}=2\Z f_2\oplus\Z f_3\oplus2\Z f_{-2}
 \cong U(4)\oplus\angles{2}.
\]
Here $U$ is the unimodular hyperbolic plane, and $U(4)$ has its
pairing multiplied by four.
The Gram-index map is
\[
 \alpha:\Gammaind=\Z^3\longrightarrow L_{II},
 \qquad
 \alpha(n,\ell,m)=2nf_2-\ell f_3+2mf_{-2}.
\]
Then
\[
 (\alpha(n,\ell,m),\alpha(n',\ell',m'))
 =-2\bigl(2nm'+2n'm-\ell\ell'\bigr),
\]
and
\[
 (\alpha(n,\ell,m),\alpha(n,\ell,m))
 =-2(4nm-\ell^2).
\]

\subsection{Real simple roots and Weyl vector}
Put
\[
 \delta_1=2f_2-f_3,
 \qquad
 \delta_2=2f_{-2}-f_3,
 \qquad
 \delta_3=f_3.
\]

\begin{proposition}[Type-II Cartan matrix]
The Gram matrix of the real simple roots is
\[
 \bigl((\delta_i,\delta_j)\bigr)_{i,j=1}^3
 =\begin{pmatrix}
 2&-2&-2\\
 -2&2&-2\\
 -2&-2&2
 \end{pmatrix}.
\]
The reflections
\[
 s_i(x)=x-(x,\delta_i)\delta_i
\]
are integral involutions preserving $L_{II}$ and its form.
\end{proposition}
\begin{proof}
Substitution into the basis pairing gives the matrix. Since $(\delta_i,\delta_i)=2$, the reflection formula has integral coefficients. Direct expansion gives $s_i^2=1$ and $(s_ix,s_iy)=(x,y)$.
\end{proof}

The Weyl vector is
\[
 \rho=f_2-\frac12f_3+f_{-2},
 \qquad (\rho,\delta_i)=-1.
\]
Use formal root-torus monomials
\[
 e^{-\alpha(n,\ell,m)}=Q^nR^\ell S^m,\qquad
 e^{-\rho}=Q^{1/2}R^{1/2}S^{1/2}.
\]
The substitution relating this torus to the denominator period
coordinate will be specified below.


\section{The root reflection matrices}
In the basis $(f_2,f_3,f_{-2})$, the form is
\[
 G=\begin{pmatrix}0&0&-1\\0&2&0\\-1&0&0\end{pmatrix}.
\]
The three reflection matrices are
\[
 S_1=\begin{pmatrix}1&4&4\\0&-1&-2\\0&0&1\end{pmatrix},\quad
 S_2=\begin{pmatrix}1&0&0\\-2&-1&0\\4&4&1\end{pmatrix},\quad
 S_3=\begin{pmatrix}1&0&0\\0&-1&0\\0&0&1\end{pmatrix}.
\]
Direct multiplication gives
\[
 S_i^2=I,
 \qquad S_i^{\mathsf T}GS_i=G.
\]


\section{The protected Jacobi input}
\subsection{The K3 elliptic genus}
Write \(q=e^{2\pi i\tau}\) and \(r=e^{2\pi iz}\).
Let $S$ be a K3 surface. The right-moving supersymmetric index of a unitary $\mathcal N=(4,4)$ sigma model is the elliptic genus
\[
 Z_{K3}(\tau,z)=2\phiweak(\tau,z).
\]
The normalized weak Jacobi form has weight zero, index one, and Fourier expansion
\[
 \phiweak(\tau,z)=\sum_{n\ge0}\sum_{\ell\in\Z}f(n,\ell)q^nr^\ell.
\]
In Eichler--Zagier normalization,
\[
 \phiweak(\tau,z)=4\sum_{a=2}^{4}
 \frac{\vartheta_a(\tau,z)^2}{\vartheta_a(\tau,0)^2}
 \in J^{\mathrm w}_{0,1}
\]
\cite{EZ}. The coefficient depends on the discriminant $4n-\ell^2$ and on $\ell$ modulo $2$.

\begin{theorem}[Exact low Fourier coefficients]
\label{comp:thm:jacobi-coefficients}
The first four $q$-rows are
\[
\begin{aligned}
\phiweak={}&r^{-1}+10+r\\
&+q\bigl(10r^{-2}-64r^{-1}+108-64r+10r^2\bigr)\\
&+q^2\bigl(r^{-3}+108r^{-2}-513r^{-1}+808-513r+108r^2+r^3\bigr)\\
&+q^3\bigl(-64r^{-3}+808r^{-2}-2752r^{-1}+4016\\
&\hspace{7.5em}{}-2752r+808r^2-64r^3\bigr)+O(q^4).
\end{aligned}
\]
In particular,
\[
 f(0,0)=10,
 \qquad \frac{f(0,0)}2=5.
\]
\end{theorem}
\begin{proof}
Insert the theta-series expansions into the preceding quotient formula. Work with $t=q^{1/8}$ and $x=r^{1/2}$; each quotient is an exact formal power series in $t$ with Laurent-polynomial coefficients in $x$. The fractional powers cancel in the sum. Division through $t^{24}$ gives the displayed rows. The arithmetic calculation is reproduced in Section~\ref{app:theta-computation}.
\end{proof}

\subsection{The virtual half-index}
For every $(n,\ell)$ choose the minimal super vector space
$U_{n,\ell}=\mathbb C^{\max(f(n,\ell),0)\mid\max(-f(n,\ell),0)}$. It satisfies
\[
 \sdim U_{n,\ell}=f(n,\ell).
\]
The integer \(f(n,\ell)\) is its superdimension. In the ordinary
Grothendieck group of finite-dimensional super vector spaces, the
even and odd dimensions remain separate. The signed integer is their
image under superdimension. The determinant depends only on that image. Adding an even--odd pair
preserves the determinant, but changes the vector space. The minimal
choice makes every zero coefficient contribute the zero space.

\begin{lemma}[Fock character and Berezinian]
Let $V=\bigoplus_{\gamma}V_\gamma$ be a super vector space graded
by a pointed positive lattice cone, with a proper positive height and
finite-dimensional homogeneous pieces and $V_0=0$. In the total-height completion,
\[
 \sch\bigl(\Sym(V_{\bar0})\tensor\bigwedge(V_{\bar1})\bigr)
 =\prod_\gamma(1-x^\gamma)^{-\sdim V_\gamma},
\]
and
\[
 \sdet_V(1-x)=\prod_\gamma(1-x^\gamma)^{\sdim V_\gamma}.
\]
\end{lemma}
\begin{proof}
An even one-dimensional summand contributes the geometric series $(1-x^\gamma)^{-1}$. An odd one-dimensional summand contributes $1-x^\gamma$. Multiplication over a homogeneous basis proves both identities.
\end{proof}



\section{Exact theta-series calculation}\label{app:theta-computation}
Set $t=q^{1/8}$ and $x=r^{1/2}$. Then
\[
\begin{aligned}
\vartheta_2(\tau,z)&=\sum_{n\in\Z}t^{4n^2+4n+1}x^{2n+1},\\
\vartheta_3(\tau,z)&=\sum_{n\in\Z}t^{4n^2}x^{2n},\\
\vartheta_4(\tau,z)&=\sum_{n\in\Z}(-1)^nt^{4n^2}x^{2n}.
\end{aligned}
\]
Cancel the common lowest power of $t$ in each quotient before
formal division in $\mathbb Q[x^{\pm1}][[t]]$. At $q=0$, the
second theta quotient contributes $(r+2+r^{-1})/4$ and the other
two contribute one each. Multiplication by four gives
$r^{-1}+10+r$. Successive convolution and inversion through
$t^{24}$ give the next three rows of
Theorem~\ref{comp:thm:jacobi-coefficients}. Fractional powers
cancel before conversion back to $(q,r)$.


\section{The Borcherds product}
\subsection{The effective chamber}
For \(Z=\left(\begin{smallmatrix}\tau&z\\z&\sigma\end{smallmatrix}\right)\in\mathbb H_2\),
put \(s=e^{2\pi i\sigma}\). Define
\[
\begin{aligned}
\Gammaeff={}&\set{(n,\ell,m)\in\Z^3:m>0,\ n\ge0}\\
&\cup\set{(n,\ell,0):n>0}
 \cup\set{(0,\ell,0):\ell<0}.
\end{aligned}
\]
The ordering is lexicographic in $(m,n,-\ell)$ and is precisely the ordering selected by the cusp chamber.
For $\gamma=(n,\ell,m)$ put
\[
 V_\gamma=U_{nm,\ell},
 \qquad \sdim V_\gamma=f(nm,\ell).
\]

\begin{definition}[Monic virtual determinant]
The monic determinant of the half-index is
\[
 \Dmon(Z)=q^{1/2}r^{1/2}s^{1/2}
 \sdet_V(1-x)
 =q^{1/2}r^{1/2}s^{1/2}
 \prod_{\Gammaeff}(1-q^nr^\ell s^m)^{f(nm,\ell)}.
\]
\end{definition}

\begin{theorem}[Borcherds--Gritsenko--Nikulin product]
The preceding product converges normally sufficiently far inside the chosen cusp chamber and extends to a holomorphic genus-two modular form of weight five with a character of order two. It is the monic Igusa form $\Dmon$. The theta-normalized form is
\[
 \Dtheta=2^6\Dmon.
\]
\end{theorem}
\begin{proof}
The singular theta lift of the vector-valued form corresponding to $\phiweak$ produces the product, divisor, weight, and automorphic character. The weight is $f(0,0)/2=5$. The genus-two identification with the Igusa form and the product formula are the Gritsenko--Nikulin theorems \cite{Borcherds95,Borcherds98,GN97,GN98b}. The product of the ten even theta constants has coefficient $2^6$ at $q^{1/2}r^{1/2}s^{1/2}$, fixing the scalar.
\end{proof}

\subsection{The theta constant normalization}
For \(a,b\in\{0,1\}^2\), fix the series
\[
 \vartheta[a,b](Z,0)=\sum_{u\in\mathbb Z^2}
 \exp\!\left(\pi i\bigl((u+a/2)^tZ(u+a/2)+b^tu\bigr)\right).
\]
Let \(\mathfrak E\) be the ten pairs with \(a^tb\) even, and set
\[
 \Dtheta(Z)=\prod_{(a,b)\in\mathfrak E}\vartheta[a,b](Z,0).
\]
This is the Gritsenko--Nikulin phase convention. Replacing the last
term by \(b^t(u+a/2)\) multiplies the ten-factor product by \(-1\):
only \(a=b=(1,1)\) contributes a nontrivial product phase.

\begin{proposition}[Ten-factor leading coefficient]
In the stated convention,
\[
 [q^{1/2}s^{1/2}]\Dtheta=64(r^{1/2}-r^{-1/2}),
 \qquad [q^{1/2}r^{1/2}s^{1/2}]\Dtheta=64.
\]
\end{proposition}
\begin{proof}
The four factors with \(a=(0,0)\) start with \(1\). The two
with \(a=(1,0)\) start with \(2q^{1/8}\), and the two with
\(a=(0,1)\) start with \(2s^{1/8}\). The remaining pair has
leading terms
\[
 2q^{1/8}s^{1/8}(r^{1/4}+r^{-1/4}),\qquad
 2q^{1/8}s^{1/8}(r^{1/4}-r^{-1/4}).
\]
Their product retains both elliptic monomials. Multiplication by
the other eight factors gives the formula. In particular the
normalization \(\Dtheta=64\Dmon\) fixes one coefficient and
does not discard the second term of the boundary coefficient.
\end{proof}

\begin{proposition}[Boundary factor of the Borcherds product]
The $m=0$ part of the monic product is
\[
 \prod_{n>0,\ell\in\Z}(1-q^nr^\ell)^{f(0,\ell)}
 \prod_{\ell<0}(1-r^\ell)^{f(0,\ell)}
 =(1-r^{-1})\prod_{n\ge1}
 (1-q^nr^{-1})(1-q^n)^{10}(1-q^nr).
\]
Consequently the first Fourier--Jacobi coefficient is
\[
 q^{1/2}r^{1/2}(1-r^{-1})
 \prod_{n\ge1}(1-q^nr^{-1})(1-q^n)^{10}(1-q^nr)
 =\eta(\tau)^9\vartheta_{11}(\tau,z),
\]
where the odd theta function is fixed by the exact convention
\[
 \vartheta_{11}(\tau,z)
 =-q^{1/8}r^{-1/2}
 \prod_{n\ge1}(1-q^{n-1}r)(1-q^nr^{-1})(1-q^n).
\]
\end{proposition}
\begin{proof}
The $q^0$ row of $\phiweak$ has exactly three nonzero coefficients:
$f(0,-1)=1$, $f(0,0)=10$, and $f(0,1)=1$. Substitution into the $m=0$ factors gives the first formula. The Jacobi triple product and $\eta(\tau)=q^{1/24}\prod_{n\ge1}(1-q^n)$ give the second, including its leading sign and monomial.
\end{proof}

\begin{corollary}[The scalar square]
The monic weight-ten form is
\[
 \ChiTen=\Dmon^2=2^{-12}\Dtheta^2.
\]
Its automorphic character is trivial.
\end{corollary}
\begin{proof}
The identity $\Dtheta=2^6\Dmon$ gives the scalar factor.  The character of $\Dmon$ has order two, so its square is trivial.
\end{proof}

\section{Divisor, character, and the meaning of the square root}
Let \(H_1\) be the image of \(z=0\) in the genus-two modular
stack. Its inverse image in \(\mathbb H_2\) is the orbit of
that hypersurface.

\begin{theorem}[Diagonal divisor]
On \(\mathbb H_2\), the divisor of $\Dmon$ is the orbit of
\(z=0\), with multiplicity one. Equivalently its divisor on the
modular stack is \(H_1\). The form transforms by the Maass character $\nu_5$ of order two, and $\ChiTen=\Dmon^2$ is scalar-valued.
\end{theorem}
\begin{proof}
This is the defining diagonal-divisor property of the Igusa square root and the first row of the Cl\'ery--Gritsenko classification \cite{Igusa62,Igusa64,GC}.
\end{proof}

\begin{proposition}[Linear and slash reflection signs]
\label{comp:prop:reflection-signs}
For the real reflections, the linear tube action has
\[
 \Dmon(s_{\delta_i}Z)=-\Dmon(Z).
\]
For the block-diagonal symplectic lifts
\(g_i=\operatorname{diag}(U_i^{-t},U_i)\), with
\[
 U_1=\begin{pmatrix}1&-2\\0&-1\end{pmatrix},\quad
 U_2=\begin{pmatrix}1&0\\2&-1\end{pmatrix},\quad
 U_3=\begin{pmatrix}1&0\\0&-1\end{pmatrix},
\]
the Maass character satisfies \(\nu_5(g_i)=1\). Thus the
homogeneous linear reflection character is the sign character of
\(W\), while these weight-five slash actions have value \(+1\).
\end{proposition}
\begin{proof}
The maps \(Z\mapsto U_i^{-t}ZU_i^{-1}\) give the three tube
reflections. Every \(U_i\) has determinant \(-1\), and each
\(g_i\) reduces to the identity modulo two. The Maass character
is the sign of the permutation of the six odd characteristics,
so \(\nu_5(g_i)=1\). The automorphy factor is
\(j(g_i,Z)=\det U_i=-1\). Hence
\(\Dmon(g_iZ)=\nu_5(g_i)j(g_i,Z)^5\Dmon(Z)=-\Dmon(Z)\).
Products of these reflections give the sign character in the linear
action. Dividing by \(j^5\) gives the stated slash action.
\end{proof}

Real roots of the denominator algebra define Weyl reflections. Imaginary roots contribute root spaces and product exponents. This division of labor is intrinsic to a Borcherds algebra: the Weyl group is generated by real simple roots, while the imaginary simple roots enter the correction term of the denominator identity.

\subsection{Automorphic section, orientation line, and Pfaffian operator}
An automorphic square root is a section of a line bundle with character. An orientation is a square root of a determinant line. A Pfaffian operator is a skew family whose regularized Pfaffian is a section of that orientation line. These form a sequence of successive refinements:
\[
\begin{gathered}
 \text{automorphic section}
 \longrightarrow \text{determinant-line identification}\\
 \longrightarrow \text{oriented square-root line}
 \longrightarrow \text{Fredholm family and Pfaffian section}.
\end{gathered}
\]

A geometric comparison must therefore provide its determinant line,
a coherent square root, and a section-preserving map to the modular
line. An operator model also specifies its domain and regularization;
a differential-operator model specifies its symbol and analytic
conditions. Chapter~\ref{comp:ch:root-operators} constructs a skew
Fredholm model of the product and retains the global geometric
comparison as a separate problem.

\section{The minimal simple-root presentation}
The real roots form a basis of \(L_{II}\). Write
\[
 \alpha=b_1\delta_1+b_2\delta_2+b_3\delta_3,
 \qquad h(\alpha)=b_1+b_2+b_3.
\]
For \(\alpha=\alpha(n,\ell,m)\), this is
\[
 (b_1,b_2,b_3)=(n,m,n+m-\ell),\qquad
 h(\alpha)=2n+2m-\ell.
\]
The active product support lies in \(Q_+=\sum_i\mathbb N\delta_i\).
Indeed, weak index-one Jacobi coefficients vanish when
\(4nm-\ell^2<-1\). For \(n,m\geq0\) with \(n+m>0\), this implies
\(\ell\leq n+m\), except possibly \(n=m\), where
integrality gives the same inequality. The boundary factor has
\((n,\ell,m)=(0,-1,0)\). Every bounded-height region is finite.

Let \(W\) be the real reflection group and put
\[
 \mathcal C=\{a=\textstyle\sum b_i\delta_i:
       b_i\geq0,\ (a,\delta_i)\leq0\text{ for all }i\}.
\]
Its generating rays are
\(a_{12}=\delta_1+\delta_2\),
\(a_{13}=\delta_1+\delta_3\), and
\(a_{23}=\delta_2+\delta_3\). To see this, write
\[
 a=t_{12}a_{12}+t_{13}a_{13}+t_{23}a_{23},\qquad
 t_{ij}=\tfrac12(b_i+b_j-b_k).
\]
The chamber inequalities say exactly \(t_{ij}\geq0\), and
\[
 (a,a)=-8(t_{12}t_{13}+t_{12}t_{23}+t_{13}t_{23}).
\]
Thus the only null rays are the three displayed ones; all other
nonzero chamber vectors are timelike.

In root-torus notation the chamber Fourier expansion has the form
\[
 \Dmon(Z_{\mathrm{root}})=\sum_{w\in W}\det(w)e^{-w\rho}
       \left(1-\sum_{0\ne a\in\mathcal C\cap L_{II}}
                         m(a)e^{-wa}\right).
\]
Here \(m(a)=-c_{\rho+a}\), where \(c_{\rho+a}\) is the
coefficient of that monomial in the monic theta product. The
reflection law makes the Fourier coefficients alternating; wall
coefficients vanish, and each remaining orbit has a unique chamber
representative. This gives the displayed expansion
\cite[Theorem~2.3]{GN97}. Along each primitive null ray,
\[
 1-\sum_{n\geq1}m(na)t^n=\prod_{n\geq1}(1-t^n)^9.
\]
The exponent nine follows from the boundary input
\(\eta^9\vartheta_{11}\), with the other rays obtained by
chamber permutations.

\begin{definition}[Minimal imaginary generators]
At a timelike chamber vector \(a\), take
\(\max(m(a),0)\) even and \(\max(-m(a),0)\) odd generators.
At every nonzero multiple of each primitive null ray, take nine
even generators and no odd generator. Together with the three
even real generators, these form a labelled set \(I_\Delta\),
with degree \(\alpha_i\) and parity \(p(i)\).
\end{definition}

The null-ray Euler exponents determine its simple multiplicities;
the coefficients \(m(na)\) are not themselves those multiplicities.
At a timelike degree the signed simple multiplicity likewise need
not equal the signed dimension of the whole root space, which also
contains brackets.

\begin{definition}[The Igusa Lie superalgebra]
Take the even Cartan space \(\mathfrak h=L_{II}\otimes\mathbb C\)
and generators \(e_i,f_i\) of parity \(p(i)\) and degrees
\(\alpha_i,-\alpha_i\). Identify \(h_i=\alpha_i\in\mathfrak h\)
and impose
\[
 [h,h']=0,\quad [h,e_i]=(h,\alpha_i)e_i,\quad
 [h,f_i]=-(h,\alpha_i)f_i,\quad [e_i,f_j]=\delta_{ij}h_i.
\]
For real \(i\ne j\), impose
\[
 (\operatorname{ad}e_i)^{1-(\alpha_i,\alpha_j)}e_j=0,
 \qquad(\operatorname{ad}f_i)^{1-(\alpha_i,\alpha_j)}f_j=0.
\]
For \((\alpha_i,\alpha_j)=0\), impose
\([e_i,e_j]=[f_i,f_j]=0\). These are super brackets, with
\([x,y]=-(-1)^{|x||y|}[y,x]\). Denote this presentation by
\(\gIg\), its positive half generated by the $e_i$ by \(\nIg\),
and its negative half generated by the $f_i$ by \(\nIg^-\).
\end{definition}

\begin{theorem}[Root supermultiplicities and denominator]
This presentation has triangular decomposition
\(\gIg=\nIg^-\oplus\mathfrak h\oplus\nIg\), with finite
even and odd root dimensions. On the active positive support,
\[
 \operatorname{sdim}\gIg_{\alpha(n,\ell,m)}=f(nm,\ell).
\]
In the Gritsenko--Nikulin period convention,
\[
 \den(\gIg)=\Dmon(2Z)=2^{-6}\Dtheta(2Z).
\]
\end{theorem}
\begin{proof}
Distinct simple vectors have nonpositive pairing: the real pairings
are \(-2\), the mixed pairings are the chamber inequalities,
and two future nonspacelike vectors pair nonpositively. Thus the
symmetrizable Borcherds presentation theorem applies. Its free
labelled Cartan reduces by the central kernel of
\(h_i\mapsto\alpha_i\), giving the stated Cartan and triangular
decomposition. In a bounded positive degree only finitely many
labelled generators and tensor words can occur, so PBW gives finite
root dimensions.

The Weyl correction of this minimal superdatum is the chamber
Fourier correction above. For the null rays the pairwise orthogonal
even generators supply the Euler product of exponent nine. Timelike
generators supply their signed coefficients. The Gritsenko--Nikulin
denominator theorem therefore identifies the root product with the
monic Igusa product \cite{GN97,GN98a,GN98b}. The change to its
period coordinate is \(Z_{\mathrm{root}}=2Z\).
\end{proof}

For example, at \(\delta_1+\delta_2+\delta_3\) there are
29 even and 93 odd vectors. Its even part consists of the two
real multilinear brackets and 27 brackets between a null-ray
simple generator and the remaining real generator. Its 93 odd
vectors are the timelike simple generators, since the correction
coefficient there is \(-93\). Thus the signed dimension is
\(-64\), while the total dimension is 122.

\chapter{Positive completions and current constructions}
\label{comp:ch:root-completions}

Positive height controls products and their limits. It gives a
finite calculation in each root degree, but distinguishes completion
of an enveloping algebra from the enveloping algebras of Lie
quotients. The ordinary bar comparison begins in that setting;
current fields and Ran envelopes then add their coefficient categories.

\section{Positive height and the ordinary bar}
\label{comp:sec:positive-bar}
Let \(Q_+=\mathbb N\delta_1+\mathbb N\delta_2+
\mathbb N\delta_3\), with height equal to the sum of the three
coefficients. Every bounded-height set of labelled simple generators
is finite. The free tensor algebra on these generators therefore has
finite-dimensional homogeneous pieces, as do its quotient
\(U(\nIg)\) and the embedded Lie superalgebra \(\nIg\).

For a positively root-graded algebra \(A\), write
\[
 A_{\leq H}=A\big/\!\bigoplus_{h(\gamma)>H}A_\gamma,
 \qquad
 \widehat A_h=\varprojlim_H A_{\leq H}
             =\prod_{\gamma\in Q_+}A_\gamma.
\]
The ideal is defined by total root height, including the height of a
product. The completed tensor product uses the same total height in
both factors. For \(A=U(\nIg)\), the primitive coproduct
\(x\mapsto x\otimes1+1\otimes x\) extends continuously to that
tensor product. At finite height its target is the quotient of
$A\otimes A$ by total height, not the untruncated tensor product
of two height quotients.

There is also a finite-dimensional nilpotent Lie quotient
\[
 \mathfrak n_H=\nIg\big/\!\bigoplus_{h(\alpha)>H}\gIg_\alpha.
\]
It is useful for calculations in degrees at most \(H\), but
\(U(\mathfrak n_H)\) still contains arbitrarily long products
of low-height generators. Thus \(\varprojlim_HU(\mathfrak n_H)\)
and \(\widehat U_h(\nIg)\) are different constructions.
For the one-dimensional even abelian Lie algebra \(\mathbb Cx\)
of height one, the former limit is \(\mathbb C[x]\), whereas
the latter is \(\mathbb C[[x]]\).

Let \(B_{\mathrm{word}}A\) be the normalized coaugmented
ordinary bar coalgebra of an augmented algebra $A$ with $A_0=\mathbb C$.
Write $\overline A$ for its augmentation ideal. Its underlying graded object is
\(T^c(s\overline A)\), with deconcatenation and the differential
from adjacent multiplication. Here \(s\) raises homological
degree by one; intrinsic parity is recorded separately.
Complete this coalgebra by the total height of its inputs, and use
the same completion for the Chevalley chains
\(C_*(\nIg)=\operatorname{Sym}(s\nIg)\).

\begin{theorem}[Positive-height bar--Chevalley comparison]
\label{comp:thm:positive-bar}
Suspended antisymmetrization gives a natural quasi-isomorphism
of height-complete dg coalgebras
\[
 \widehat C_{*,h}(\nIg)\longrightarrow
 \widehat B_{\mathrm{word},h}U(\nIg).
\]
In every fixed root degree it is already obtained from any
\(\mathfrak n_H\) containing that degree. Its Euler
supercharacter is
\[
 \chi_{\mathrm{root}}\widehat C_{*,h}(\nIg)
 =\prod_{\alpha>0}(1-e^{-\alpha})^{\operatorname{sdim}\gIg_\alpha}.
\]
\end{theorem}
\begin{proof}
The map is
\[
 sx_1\cdots sx_r\longmapsto
 \sum_{\sigma\in S_r}\varepsilon(\sigma;sx_1,\ldots,sx_r)
 [x_{\sigma(1)}|\cdots|x_{\sigma(r)}].
\]
Here $\varepsilon$ is the permutation sign in the symmetry of
complexes of super vector spaces: homological degree and intrinsic
parity each contribute their Koszul sign. There is no factor \(1/r!\): the unshuffle coproduct on the
source agrees with deconcatenation of the permutation sum. The
binary differential agrees because the commutator in the enveloping
algebra is the given bracket. PBW reduces the filtered comparison
to the Koszul resolution of the super-symmetric algebra.
In each fixed root degree all inputs have bounded positive height,
and the filtration and chain groups are finite. The ordinary
Koszul contraction therefore proves the quasi-isomorphism there.
Products of complexes of vector spaces preserve these
quasi-isomorphisms, giving the completed map.

An intrinsically even root vector becomes an odd Chevalley generator
and contributes \(1-e^{-\alpha}\) to the Euler supercharacter.
An intrinsically odd root vector contributes its reciprocal.
The parity sign and the homological Euler sign together give the
displayed exponent. Positive height makes each coefficient finite.
\end{proof}

\begin{proposition}[The bar of the Lie-quotient tower]
\label{comp:prop:lie-quotient-bar}
There is also a natural quasi-isomorphism of inverse-limit complexes
\[
 \varprojlim_H C_*(\mathfrak n_H)
       \longrightarrow
 \varprojlim_H B_{\mathrm{word}}U(\mathfrak n_H).
\]
In each homological degree the left side is
\(\widehat C_{*,h}(\nIg)\). The natural coefficient map from
the right side to \(\widehat B_{\mathrm{word},h}U(\nIg)\)
is a quasi-isomorphism as well. These comparisons do not identify
\(\varprojlim_HU(\mathfrak n_H)\) with \(\widehat U_h(\nIg)\).
\end{proposition}
\begin{proof}
At each \(H\), ordinary antisymmetrization is a quasi-isomorphism
by PBW, and the comparisons commute with the quotient maps. The
bonding maps on Chevalley and bar chain groups are surjective:
they are induced by surjections of Lie algebras, enveloping
algebras and augmentation ideals. The same is true on the mapping
cones \(K_H\) of the comparisons. Consequently the Roos sequence
\[
 0\longrightarrow\varprojlim_HK_H\longrightarrow
 \prod_HK_H\xrightarrow{\,1-\mathrm{shift}\,}\prod_HK_H
 \longrightarrow0
\]
is degreewise exact. Products of complexes of vector spaces preserve
acyclicity. Since every \(K_H\) is acyclic, its inverse limit
is acyclic, proving the first assertion.

In Chevalley degree \(r\), a monomial with inputs of height at
most \(H\) has total height at most \(rH\). Conversely, a
monomial of total height at most \(H\) has all inputs of height
at most \(H\). These two filtrations are therefore cofinal in
each chain degree, identifying the completed Chevalley complexes.
Every fixed root coefficient of the bar tower also stabilizes,
giving the stated coefficient map. Its composite with
antisymmetrization is the map of
Theorem~\ref{comp:thm:positive-bar}, so it is a quasi-isomorphism.
The example \(\mathbb C[x]\) versus \(\mathbb C[[x]]\) above
still distinguishes the two enveloping constructions.
\end{proof}

Write \(Z_{\mathrm{root}}\) for the root-torus coordinate with
\[
 e^{-\alpha(n,\ell,m)}=Q^nR^\ell S^m,
 \qquad e^{-\rho}=Q^{1/2}R^{1/2}S^{1/2}.
\]
The root-torus product is \(\Dmon(Z_{\mathrm{root}})\).
The Gritsenko--Nikulin period convention substitutes
\(Z_{\mathrm{root}}=2Z\). Consequently
\[
 e^{-\rho}\chi_{\mathrm{root}}\widehat C_{*,h}(\nIg)
 =\den(\gIg)=\Dmon(2Z).
\]
This is an ordinary word-bar identity. An associative algebra on
Ran space and a chiral field realization require their own
coefficient categories and comparison maps.

\begin{proposition}[The two-sided root product]
\label{comp:prop:two-sided-product}
The Lie bracket has no jointly continuous extension to the
coordinatewise product of all nonzero root spaces with discrete
Cartan factor.
\end{proposition}
\begin{proof}
Choose one even simple pair \(e_n,f_n\) at each degree
\(na,-na\) on a primitive chamber null ray. The presentation gives
\([e_n,f_m]=\delta_{nm}na\). Both sequences of partial sums
converge coordinatewise, but their brackets have Cartan component
\(N(N+1)a/2\), which is not eventually constant. Continuity
would force convergence in the discrete Cartan space.
\end{proof}

The positive completion with a discrete finite-support negative
half does have a continuous bracket. Every mixed output coefficient
then receives only finitely many terms. A completion of a full
quantum double needs an analogous specified domain; root-degree
preservation alone does not provide one.
\section{Currents and their associative recipient}
\label{comp:sec:currents}
On a coordinate disc, the level-zero current Lie conformal
superalgebra is
\[
 \operatorname{Cur}(\mathfrak n_H)=\mathbb C[\partial]\otimes
 \mathfrak n_H,\qquad [a_\lambda b]=[a,b].
\]
Its coordinate-free input is the induced right differential module
\(N_H=\mathfrak n_H\otimes D_E\), with its Lie-* bracket.
Let \(V_H=U^{\mathrm{ch}}(N_H)\) be its BD current envelope
\cite{BD04}. In a coordinate, its vacuum module is spanned by ordered
negative modes \(a(-j)\), \(j\geq1\).

\begin{proposition}[Current character]
\label{comp:prop:current-character}
At fixed root degree and energy, the quotient system \((V_H)_H\)
stabilizes. Its coefficientwise supercharacter is
\[
 \prod_{j\geq1}\prod_{\alpha>0}
       (1-t^je^{-\alpha})^{-\operatorname{sdim}\gIg_\alpha}.
\]
The current generators are primitive for the level-zero current
coproduct. The coefficientwise completion does not by itself assert
uniform pole bounds for all completed states.
\end{proposition}
\begin{proof}
PBW identifies the vacuum space with the super-symmetric algebra
on \(\mathfrak n_H\otimes t^{-1}\mathbb C[t^{-1}]\).
The mode \(a(-j)\) has energy \(j\); an even mode gives a
symmetric factor and an odd mode an exterior factor. At fixed
root height and energy only finitely many modes and products
contribute. The current coproduct sends
\(a(z)\) to \(a(z)\otimes1+1\otimes a(z)\), and the
quotient maps preserve it. Lower truncation for arbitrary completed
states is a separate continuity condition on fields.
\end{proof}

An associative bar needs an associative algebra and its tensor.
Use complexes of super vector spaces, retaining intrinsic parity
separately from cochain degree. For nonempty finite sets and
surjections, form
\[
 \mathcal R_E=\varprojlim_I D\operatorname{-Mod}(E^I;\mathrm{sVect}_{\mathbb C}),
 \qquad \mathcal R_E^+=D(\mathrm{sVect}_{\mathbb C})\times\mathcal R_E,
\]
using exceptional pullback along the diagonals. For an ordered
partition \(I=I_1\sqcup I_2\), let \(j_{I_1,I_2}\) be the
locus where the two blocks are disjoint. The chiral convolution is
\[
 (F\star G)_I=\bigoplus_{I=I_1\sqcup I_2}
 j_{I_1,I_2,*}j_{I_1,I_2}^{!}(F_{I_1}\boxtimes G_{I_2}).
\]
Its unit \(K\) has only the empty coefficient \(K_\varnothing
=\mathbb C\). A positive object has zero empty coefficient.
The BD vacuum is the unit map $\omega_E\to V_H$. It is
different from this empty-component unit.

The underlying chiral Lie algebra of \(V_H\) gives the positive
main-diagonal Lie object \(\mathcal L_H=\iota_*V_H\), with
\((\iota_*V_H)_I=\delta_{I,*}V_H\). Define its augmented
associative envelope \(\mathcal A_H=U_\star(\mathcal L_H)\).

\begin{proposition}[Bar of the Ran envelope]
\label{comp:prop:ran-envelope}
In \((\mathcal R_E^+,\star,K)\),
\[
 B_\star\mathcal A_H\simeq
 \operatorname{CE}_\star(\mathcal L_H),\qquad
 \Omega_\star\operatorname{CE}_\star(\mathcal L_H)
       \simeq\mathcal A_H.
\]
The Chevalley coalgebra satisfies factorization.
\end{proposition}
\begin{proof}
A positive convolution word of length \(r\) vanishes on a
configuration of cardinality less than \(r\). Bar and cobar
filtrations are therefore finite on each configuration. Flattening
a bar--cobar word gives one two-term contractible gap complex at
every internal cut; only the empty and one-input terms survive.
This proves the associative unit and counit equivalences.

Suspended symmetrization gives the envelope comparison, with the
same permutation sum as in Theorem~\ref{comp:thm:positive-bar}.
Its PBW filtration reduces to the super-symmetric Koszul complex;
the Euler contraction divided by positive polynomial degree is
valid in characteristic zero. The finite configuration filtration
proves the equivalence. Finally, the main-diagonal support of
\(\mathcal L_H\) implies factorization of its chiral Chevalley
coalgebra \cite[Theorems~5.1.1--5.2.1]{FG12}.
\end{proof}

The last cited equivalence is Lie--commutative. The associative
assertion also uses the envelope and the bar--cobar argument.
The objects \(N_H\), \(V_H\), \(\mathcal L_H\) and
\(\mathcal A_H\) have different carriers. In particular, this
proposition does not identify \(\mathcal A_H\) with the BD
field algebra \(V_H\), or construct a second independently
defined field dual. Such a comparison must preserve its specified
supported operations and units.

\section{The primitive comparison and its extensions}
Let \(\mathfrak p\) be a positive-root-graded Lie superalgebra
with finite-dimensional homogeneous pieces and a proper positive
height. A primitive realization of the Igusa presentation is a
parity-preserving homogeneous map from its labelled simple-root
spaces to \(\mathfrak p\), satisfying the defining relations.

\begin{theorem}[Universal presentation map]
\label{igusa:primitive-presentation-map}
Let \(\mathfrak p\) be a positive-root-graded Lie superalgebra
over \(\mathbb C\), with finite-dimensional homogeneous pieces.
Suppose a homogeneous parity-preserving map from every simple-root
generating space of \(\nIg\) to \(\mathfrak p\) satisfies
the defining Borcherds--Serre relations. It extends uniquely to
a graded Lie-superalgebra map
\[
 f:\nIg\longrightarrow\mathfrak p.
\]
It induces continuous maps on height completions, enveloping
algebras, Chevalley chains, and level-zero currents. The map \(f\)
is an isomorphism if and only if it is surjective and the root-graded
PBW series agree in each parity.
\end{theorem}
\begin{proof}
Extend the chosen maps to the free Lie superalgebra on the
simple-root spaces. The relation hypothesis makes the defining
ideal lie in its kernel, giving \(f\). The generators prove
uniqueness. On enveloping algebras the map sends
\(x_1\cdots x_j\) to \(f(x_1)\cdots f(x_j)\).
On Chevalley chains it is the graded symmetric map on the shifted
generators. It commutes with the Chevalley differential because
\(f\) preserves brackets. The current map is
\(\partial^j x\mapsto\partial^j f(x)\), preserving
\([x_\lambda y]=[x,y]\). Root degree preservation gives
compatible maps at each height and hence on the stated limits.

If \(f\) is surjective, \(U(f)\) is surjective. Proper
positive height makes each root-graded PBW component
finite-dimensional. Equality of its even and odd dimensions makes
\(U(f)\) bijective in every component. PBW embeds each Lie
superalgebra into its enveloping algebra, so \(f\) is injective.
The converse follows by applying the constructions to its inverse.
\end{proof}


A geometric comparison also requires functors on the actual Hall
and critical categories, a compatible coproduct and Hopf pairing,
and maps of determinant lines and operator domains. These categories
and lines are not the domain of the primitive Lie map. The
additional maps must be constructed and checked against that map.
They are not consequences of its universal property.


\section{The local--wrapped root algebra}
The positive cone is graded by $\gamma=(n,\ell,m)$, with elliptic degree $m$.  Set
\[
 \nIg_{\loc}=\bigoplus_{\substack{\alpha>0\\m(\alpha)=0}}\gIg_\alpha,
 \qquad
 \nIg_{\wrap}=\bigoplus_{\substack{\alpha>0\\m(\alpha)>0}}\gIg_\alpha.
\]

\begin{proposition}[Semidirect local--wrapped decomposition]
The space $\nIg_{\loc}$ is a Lie subsuperalgebra, $\nIg_{\wrap}$ is a Lie ideal, and
\[
 \nIg=\nIg_{\loc}\ltimes\nIg_{\wrap}.
\]
The proper height filtration preserves both summands.
\end{proposition}
\begin{proof}
The bracket adds root degrees.  The sum of two roots of elliptic degree zero again has degree zero, and addition of a nonnegative degree to a positive degree remains positive.  Height is additive on every nonzero bracket.
\end{proof}

Write \(\mathbb U_{\Delta_5}=\widehat U_h(\nIg)\), and
retain the current system \((V_H)_H\) of
Section~\ref{comp:sec:currents}. The adjoint action of
\(\nIg_{\loc}\) on \(\nIg_{\wrap}\) extends to derivations
of its enveloping algebra. Let \(\#\) denote the resulting
smash product, completed by total root height.

\begin{theorem}[Local--wrapped enveloping algebra]
There are natural isomorphisms
\[
 \widehat U_h(\nIg)
 \cong\widehat U_h(\nIg_{\wrap})\#\widehat U_h(\nIg_{\loc}),
\]
and
\[
 \widehat C_{*,h}(\nIg)
 \cong\widehat{\operatorname{Tot}}_h
 C_*(\nIg_{\loc};C_*(\nIg_{\wrap})).
\]
The primitive Lie superalgebra of the completed enveloping algebra
is \(\widehat{\nIg}_h\). The Weyl-shifted Euler supercharacter
of the completed Chevalley complex is \(\Dmon(2Z)\).
\end{theorem}
\begin{proof}
PBW orders wrapped generators before local generators. Moving a
local generator past a wrapped one adds their adjoint bracket,
which is exactly the smash-product relation. All relations preserve
total height, so the isomorphism extends to the completion.
Splitting a Chevalley word into its two types gives the
Hochschild--Serre bicomplex: its differentials are the two internal
brackets and the adjoint action. This proves the second isomorphism.

In characteristic zero, the primitives of a super-symmetric Hopf
algebra are its degree-one generators. Apply PBW to the highest
polynomial degree of a primitive in each finite root-degree piece.
All higher degrees vanish. This identifies the primitives before
and after completion. The denominator identity gives the character.
\end{proof}

These operations define the algebraic semidirect-product example
of a local--wrapped theory. Extending the local Lie action to chiral
currents acting on geometric wrapped objects requires supported
coefficient maps. The adjoint action alone does not specify those
maps.

\section{The colored elliptic Ran space}
The elliptic degree is a global charge carried by the whole configuration.  Define
\[
 \Ran^{\mathrm{wr}}(E)=\coprod_{m\ge0}\Ran(E)_m
\]
with monoidal operation
\[
 (I,m)\sqcup(J,n)=(I\sqcup J,m+n)
\]
on disjoint supports. The names local and wrapped refer here to
the degree colors $m=0$ and $m>0$; the underlying current operations
remain supported on configurations of points.

Let
\[
 \nIg=\bigoplus_{m\ge0}\nIg[m]
\]
be the elliptic-degree grading.  The root bracket obeys
\[
 [\nIg[m],\nIg[n]]\subset\nIg[m+n].
\]

\begin{construction}[Elliptic-degree factorization coalgebras]
For each \(H\), take
\[
 \mathcal F_H=\operatorname{CE}_\star(\mathcal L_H),
 \qquad \mathcal L_H=\iota_*V_H.
\]
Every current state has an elliptic degree, obtained by adding the
root degrees of its negative modes. Give each Chevalley word the
sum of these degrees and place its degree-\(m\) coefficient on
\(\operatorname{Ran}(E)_m\). The empty coefficient remains the
unit \(K\) defined in Section~\ref{comp:sec:currents}.
\end{construction}

\begin{theorem}[Elliptic-degree factorization]
The coalgebras \(\mathcal F_H\) factorize compatibly with the
elliptic-degree grading and the quotient system in \(H\).
Disjoint-union structure maps add the degree labels. Under the
Lie--commutative Koszul equivalence, their derived Lie primitives
are \(\mathcal L_H\). The energy-one current states recover
\(\mathfrak n_H\), with its bracket given by the zeroth product.
\end{theorem}
\begin{proof}
The current bracket and the Chevalley differential preserve root
degree. The factorization theorem of
Proposition~\ref{comp:prop:ran-envelope} therefore respects its
elliptic component. Its Lie inverse recovers \(\mathcal L_H\).
In a coordinate, PBW identifies the energy-one subspace of \(V_H\)
with the states \(a(-1)|0\rangle\), for \(a\in\mathfrak n_H\).
Their zeroth product is \([a,b](-1)|0\rangle\). All constructions
commute with the root-graded quotient maps.
\end{proof}

The degree-zero sector acts on every positive-degree sector through
the current operations. This is a grading of a point-supported
current construction. A compact sheaf projecting onto the whole
elliptic curve has different support. Its comparison must map the
common-torsor extension correspondences to these current operations,
preserving both support and charge.

\section{Cotangent fields and finite BV models}
Choose a holomorphic volume form \(\Omega_E\) on \(E\).
For each finite-height Lie algebra \(\mathfrak n_H\), the local
dg Lie algebra on \(E\times\mathbb R\) is
\[
 \mathfrak L_H=\Omega^{0,*}(E)\widehat\otimes
             \Omega^*(\mathbb R)\otimes\mathfrak n_H,
 \qquad d=\bar\partial+d_{\mathbb R}.
\]
The tensor product uses smooth local sections. With cohomological
shifts, the cotangent fields are
\[
 \mathcal E_H=
 \bigl(\Omega^{0,*}(E)\widehat\otimes\Omega^*(\mathbb R)\bigr)
 \otimes\bigl(\mathfrak n_H[1]\oplus\mathfrak n_H^\vee\bigr).
\]
The first summand is \(A\) and the second is \(B\).
Integration against \(\Omega_E\) makes their pairing have
degree \(-1\), so this is the odd cotangent theory of
\(\mathfrak L_H[1]\).
Its action, on compactly supported fields, is
\[
 S_H(A,B)=\int_{E\times\mathbb R}\Omega_E\wedge
 \left\langle B,dA+\tfrac12[A,A]\right\rangle.
\]

\begin{proposition}[Classical cotangent action]
This action satisfies the classical master equation. Its classical
local observables form a factorization algebra. The same construction
on \(\mathbb R^2_{\mathrm{top}}\times E\) uses
\[
 \mathfrak L_H^{(2)}=\Omega^*(\mathbb R^2)\widehat\otimes
                   \Omega^{0,*}(E)\otimes\mathfrak n_H
\]
with fields in
$\Omega^*(\mathbb R^2)\widehat\otimes\Omega^{0,*}(E)\otimes
(\mathfrak n_H[1]\oplus\mathfrak n_H^\vee[1])$
and action
\[
 S_H^{(2)}=\int_{\mathbb R^2\times E}\Omega_E\wedge
       \langle B,(d+\bar\partial)A+\tfrac12[A,A]\rangle.
\]
\end{proposition}
\begin{proof}
The vector field \(Q(A)=dA+[A,A]/2\) squares to zero because
\(d^2=0\), \(d\) is a derivation, and the bracket satisfies
Jacobi. Its cotangent lift is Hamiltonian with Hamiltonian \(S_H\).
Thus \(\{S_H,S_H\}=0\). Restriction of fields and extension of
compactly supported test functions give the local observable maps;
disjoint supports give their factorization products. The
Dolbeault--de Rham complex is elliptic; its local descent gives
Weiss descent for the classical cotangent observables
\cite{CG1,CG2}. The same argument applies with the second
topological direction.
\end{proof}

There is a separate finite-dimensional quantum calculation. Put
\(M_H=\mathfrak n_H[1]\), and take its odd cotangent space
\(T^*[-1]M_H\). In linear coordinates \(c^a\) and their
cotangent coordinates \(b_a\), let \(Q\) be the Chevalley
vector field and set
\[
 S_H^{\mathrm{fin}}=\sum_a b_aQ(c^a).
\]
Use the constant Berezin volume to define the BV Laplacian \(\Delta\).

\begin{proposition}[Finite cotangent quantum master equation]
The finite action satisfies
\[
 \tfrac12\{S_H^{\mathrm{fin}},S_H^{\mathrm{fin}}\}
       +\hbar\Delta S_H^{\mathrm{fin}}=0.
\]
\end{proposition}
\begin{proof}
The classical term vanishes by \(Q^2=0\). The BV term is the
divergence of \(Q\); its linear coefficients are the modular
supertrace \(x\mapsto\operatorname{str}(\operatorname{ad}x)\).
For an even homogeneous \(x\), the adjoint action strictly raises
positive root height and is upper triangular with zero diagonal.
For an odd \(x\), it is an odd endomorphism and has zero
supertrace. Linearity proves that the entire divergence vanishes.
\end{proof}

A reduction of the local field theory must specify a contraction of
its field complex and the induced pairing and operations. Reducing
only a topological interval leaves the Dolbeault fields on \(E\).
Even their harmonic complex contains
\(H^{0,*}(E)=\mathbb C\oplus\mathbb C\eta\), with \(\eta\)
in degree one. It is not the single copy \(\mathfrak n_H[1]\)
used in the finite proposition. A quantum pushforward must in
addition control its measure, boundary conditions and loop terms.

The algebraic boundary model can be specified independently:
quantize the linear Lie--Poisson space by \(U(\mathfrak n_H)\)
and take its two augmentation modules. Their derived tensor product is
\[
 \mathbb C\otimes^{\mathbb L}_{U(\mathfrak n_H)}\mathbb C
       \simeq C_*(\mathfrak n_H).
\]
The bar--Chevalley map proves this statement. The holomorphic current
model is \(V_H\). Identifying these algebraic objects with boundary
observables and an interval amplitude of the field theory requires
boundary quantization and sewing maps. At infinite height the
appropriate total-height completion must also be preserved; the
limit of the uncompleted \(U(\mathfrak n_H)\) is different.
Dualizing the Lie quotient maps reverses their direction on the
$B$-fields, so a field-theoretic limit also requires compatible
pushforward maps.

\chapter{Extensions, charges and quiver Hall models}
\label{comp:ch:quiver-models}

A geometric Hall product remembers extensions and their relative
elliptic position. Its first construction is a correspondence over
the stack of common translation torsors. Quiver models provide
explicit algebraic examples, while critical coefficients and their
orientation lines impose additional equations on the same flags.

\section{Extensions over the elliptic classifying stack}
Work on locally noetherian complex schemes with the fppf topology.
Let \(BE\) be the stack of right \(E\)-torsors. For a left
\(E\)-stack \(M\), an object of \([M/E](T)\) is a right
torsor \(P\to T\) and an equivariant map \(P\to M\).
Forgetting the map gives \([M/E]\to BE\). All products of
stacks below are two-fibre products.

\begin{proposition}[The common translation torsor]
\label{igusa:platonic-common-torsor}
There is a natural equivalence
\[
 [(M\times N)/E_{\mathrm{diag}}]
 \simeq[M/E]\times_{BE}[N/E].
\]
The equivalences for any finite number of factors commute with
regrouping and insertion of the terminal \(E\)-stack.
\end{proposition}
\begin{proof}
An object on the right is \((P,m),(Q,n)\) together with a
torsor isomorphism \(\alpha:P\to Q\). Send it to
\((P,m,n\alpha)\). Conversely, use a common torsor twice
and its identity isomorphism. A pair of torsor maps \(a,b\)
on the right satisfies \(b\alpha=\alpha'a\), so \(b\)
is determined by \(a\). The same construction identifies the
maps and their descent isomorphisms. It commutes with base change.
For several factors, every regrouping gives one torsor with the
same list of equivariant maps, proving all stated coherences.
\end{proof}

For \(M=N=E\) with translation, the diagonal quotient is
\(E\), with coordinate \(y-x\). Both separate quotients
are points. Their fibre product over \(BE\) retains the
relative translation, while their ordinary product loses it.

Let \(\mathfrak M\) be the stack of coherent sheaves on
\(X=S\times E\), with families flat over the parameter
scheme, and put \(\mathfrak Q=[\mathfrak M/E]\).
An object of \(\mathfrak Q(T)\) is a torsor \(P\to T\)
and a flat sheaf on \(S\times P\). Fix a polarization and
write \(\mathfrak Q_\alpha\) for a Hilbert-polynomial component.
Translation preserves these polynomials. Let
\(\mathfrak H_{\alpha,\beta}\) classify short exact sequences
\(0\to A\to F\to B\to0\) on a common \(S\times P\),
with all terms flat over \(T\) and polynomials
\(\alpha,\alpha+\beta,\beta\). Its span over \(BE\) is
\[
 \mathfrak Q_\alpha\times_{BE}\mathfrak Q_\beta
 \xleftarrow{p}\mathfrak H_{\alpha,\beta}
 \xrightarrow{q}\mathfrak Q_{\alpha+\beta}.
\]
The category of spans of fppf stacks over \(BE\) has composition
by fibre product and tensor product \(\times_{BE}\). Its unit
object is the identity map of \(BE\).

\begin{theorem}[Geometric Hall correspondence]
\label{igusa:platonic-common-torsor-hall}
For fixed \(\alpha,\beta\), \(q\) is representable by
algebraic spaces and proper. The spans form a coherently associative
unital Hall object in the category of spans over \(BE\).
Its unit is the zero-sheaf span
\(BE\leftarrow BE\to\mathfrak Q\).
\end{theorem}
\begin{proof}
Over a fixed \((P,F)\), the fibre of \(q\) classifies
quotients \(F\twoheadrightarrow B\) of polynomial \(\beta\).
Its kernel is flat, because \(F\) and \(B\) are flat.
An automorphism fixing the quotient and \(F\) is the identity.
Trivializing \(P\) fppf-locally identifies this fibre with the
fixed-polynomial Quot scheme, projective over the parameter scheme
by Grothendieck's theorem \cite[Theorem~3.2]{IgusaGrothendieckQuot1961}.
The quotient functor descends to an algebraic space, and properness
descends with it.

Both composites of three Hall spans classify
\(0\subset A\subset D\subset F\) on the same torsor.
One bracketing concatenates inclusions. The other starts with
\(A\subset F\) and a subobject \(B\subset F/A\), and sets
\(D=F\times_{F/A}B\). These constructions are inverse on
objects and isomorphisms. They preserve flatness and commute with
base change. For longer flags all regroupings identify the same
subobjects and quotients. Their comparison maps agree, proving the
pentagon and higher coherences. Finally \(\mathfrak Q_0=BE\).
The sequences with zero subobject or zero quotient give the unit
and its coherences.
\end{proof}

Properness here concerns the full coherent-sheaf stack on a fixed
pair of Hilbert-polynomial components. A stability restriction can
replace its fibre by an open part of a Quot scheme, which need not
be proper. A bounded charge window also needs closure under the
extensions and flags used in its operations. A reduced critical-coefficient realization must supply pullback,
pushforward, and orientation comparisons on these same flags. Its unit has
stabilizer \(E\). A coefficient theory for this quotient must
include that nonaffine stabilizer, or work relative to \(BE\),
on the unit and every flag.


\section{Additive physical charge and quadratic automorphic degree}
Let
\[
 \Lam=\widetilde H(S,\Z)
\]
be the Mukai lattice, with pairing
\[
 (r,c,s)\cdot(r',c',s')=c\cdot c'-rs'-r's
\]
and signature $(4,20)$. The electric--magnetic charge lattice is
\[
 \Gammphys=\Lam\oplus\Lam,
 \qquad c=(Q,P).
\]
Its addition is the physical superselection law. Define the Gram map
\[
 \PiGram(Q,P)=
 \left(\frac{Q^2}{2},Q\cdot P,\frac{P^2}{2}\right)
 \in\Gammaind.
\]

\begin{proposition}[Polarization identity]
For $c_i=(Q_i,P_i)$,
\[
 \PiGram(c_1+c_2)=\PiGram(c_1)+\PiGram(c_2)+B(c_1,c_2),
\]
where
\[
 B(c_1,c_2)=
 \bigl(Q_1\cdot Q_2,
 Q_1\cdot P_2+Q_2\cdot P_1,
 P_1\cdot P_2\bigr).
\]
The cross-term is a symmetric vector-valued two-cocycle:
\[
 B(c_1,c_2)+B(c_1+c_2,c_3)
 =B(c_2,c_3)+B(c_1,c_2+c_3).
\]
\end{proposition}
\begin{proof}
Expand the three quadratic expressions. The cocycle identity is bilinearity of $B$.
\end{proof}

\subsection{The polarized central extension}
Define the abelian central extension
\[
 0\longrightarrow\Gammaind\longrightarrow
 \widetilde\Gamma_\Pi
 \longrightarrow\Gammphys\longrightarrow0,
 \qquad
 \widetilde\Gamma_\Pi=\Gammphys\times\Gammaind,
\]
with multiplication
\[
 (c,\lambda)\cdot(d,\mu)
 =\bigl(c+d,\lambda+\mu+B(c,d)\bigr).
\]
The cocycle identity makes this multiplication associative. In fact
$(c,\lambda)\mapsto(c,\lambda-\PiGram(c))$ identifies it with
the direct-product group, so the extension is split. The
polarization identity makes the graph map
\[
 j_\Pi:\Gammphys\longrightarrow\widetilde\Gamma_\Pi,
 \qquad c\longmapsto(c,\PiGram(c))
\]
a homomorphism.

Let $A=\bigoplus_{c\in\Gammphys}A_c$ be a charge-graded associative algebra, and let $\kfield[\Gammaind]$ have monomials $u^\lambda$. Define the polarized central-extension algebra
\[
 \mathcal R_\Pi(A)=A\tensor\kfield[\Gammaind]
\]
with product
\[
 (a_cu^\lambda)\star(b_du^\mu)
 =a_cb_d\,u^{\lambda+\mu+B(c,d)}.
\]
Give $a_cu^\lambda$ degree $(c,\lambda)\in\widetilde\Gamma_\Pi$.

\begin{theorem}[Graph-multiplicative polarization]
The product $\star$ is associative and $\widetilde\Gamma_\Pi$-graded. The map
\[
 \iota_\Pi:A\longrightarrow\mathcal R_\Pi(A),
 \qquad a_c\longmapsto a_cu^{\PiGram(c)}
\]
is an algebra homomorphism whose homogeneous degree is the graph element $j_\Pi(c)$.
\end{theorem}
\begin{proof}
The two parenthesizations differ by
\[
 B(c,d)+B(c+d,e)-B(d,e)-B(c,d+e),
\]
which vanishes by the cocycle identity. The polarization formula gives
\[
 \iota_\Pi(a_c)\star\iota_\Pi(b_d)
 =a_cb_d\,u^{\PiGram(c+d)}.
\]
Thus the graph of the quadratic Gram map is multiplicative inside the central extension.
\end{proof}

The same construction applies to Lie superalgebras and their enveloping algebras. It retains additive physical charge, Gram degree, and the interaction cross-term in one grading group.

\section{Point-local and wrapped sectors}
Let $p_E:S\times E\to E$ be projection.

\begin{lemma}[Positive elliptic degree is wrapped]
Let $C\subset S\times E$ be a proper one-dimensional subscheme whose pushforward $(p_E)_*[C]$ has positive degree. Then $p_E(C)=E$.
\end{lemma}
\begin{proof}
Properness makes $p_E(C)$ closed. Positive pushforward degree implies that at least one irreducible component of $C$ dominates a one-dimensional closed subset of $E$. Every one-dimensional closed subset of the irreducible curve $E$ is $E$ itself, so the image of that component, and hence the image of $C$, is all of $E$.
\end{proof}

Objects whose support has finite image in $E$ give the point-local
sector. Positive elliptic degree gives support over the whole curve.

\begin{definition}[Mixed geometric operations]
Fix a symmetric monoidal coefficient theory on the relevant
configuration, object and flag stacks, with its pullbacks and
pushforwards. A local--wrapped datum consists of a local
factorization object \(\mathcal F_{\loc}\), wrapped coefficient
objects \(W_m\) for \(m>0\), ordered products
\[
 W_{m_1}\otimes\cdots\otimes W_{m_r}
       \longrightarrow W_{m_1+\cdots+m_r},
\]
and mixed maps from local insertions and wrapped inputs to wrapped
outputs on their specified incidence diagrams. These maps must
commute with disjoint local factorization, regrouping of wrapped
extensions, substitution of mixed operations, and insertion of the
unit. A geometric realization supplies the diagrams and coefficient
maps, not only the resulting grading.
\end{definition}

Write $\mathsf{Ch}_E(k)$ for a chosen local chiral operation object
and $\mathsf{Ass}(r)$ for the ordered associative operation object.
Mixed substitution is additional data. The expression
\(\mathsf{Ch}_E(k)\otimes\mathsf{Ass}(r)\) alone does not define
it: substituting several wrapped products requires a rule assigning
the local insertions to their factors. One can form the free colored
operad on local, wrapped and mixed generators and then impose the
identities of a supplied geometric datum. Identifying its operation
objects with the simple tensor expression requires a proved
normal-form theorem and the required distribution maps.

There is a concrete algebraic example. Take
\(L=\nIg_{\loc}\), \(W=U(\nIg_{\wrap})\), and let \(L\)
act by adjoint derivations. For homogeneous \(x\in L\) and
\(u,v\in W\),
\[
 x\cdot(uv)=(x\cdot u)v+(-1)^{|x||u|}u(x\cdot v),
 \qquad x\cdot1=0.
\]
Jacobi makes this a Lie action. It extends to a module-algebra
action of \(U(L)\), whose primitive coproduct is the distribution
rule in the displayed formula. Thus the colored operad generated
by a Lie bracket, an associative product and a derivation action has
this local--wrapped algebra. The total-height completion preserves
it. Chiral currents on a curve and compact wrapped coefficients
need their supported analogues of these maps.

For \(S\times E\), common-torsor extension and mixed-support
flag correspondences provide the geometric diagrams to compare.
A critical coefficient system satisfying base change and the mixed
substitution equations would realize the preceding datum. Properness
of a middle-term map alone supplies none of its critical pullback
or collision maps.

\subsection{Two product conventions for the Gram graph}
The positive cocycle in \(\mathcal R_\Pi(A)\) makes
\(a_c\mapsto a_cu^{\PiGram(c)}\) a homomorphism. In the
ordinary untwisted product on \(A\otimes\mathbb C[\Gammaind]\),
the same graph monomials instead satisfy
\[
 (a_cu^{\PiGram(c)})(b_du^{\PiGram(d)})
 =u^{-B(c,d)}\,a_cb_du^{\PiGram(c+d)}.
\]
This follows by subtracting the polarization identity. The negative
cocycle therefore describes multiplication in graph coordinates
for the untwisted algebra. The positive cocycle describes the
changed multiplication that makes the graph map multiplicative.
Both retain additive physical charge; neither makes the quadratic
Gram map additive in its original target \(\Gammaind\).

\section{Finite simple-root windows}
Let $I_\Delta$ be the multiset of simple roots of the Gritsenko--Nikulin Borcherds superalgebra.  A repeated imaginary simple root is written once for each copy.  The parity map is
\[
 p:I_\Delta\longrightarrow\Z/2.
\]
Use the positive root height \(h(i)=h(\alpha_i)\), whose
sublevel sets of labelled generators are finite.  Put
\[
 I_H=\set{i\in I_\Delta:h(i)\le H}.
\]
The restriction of the invariant form gives a finite symmetric Borcherds--Cartan matrix
\[
\begin{aligned}
 A_H&=(a_{ij})_{i,j\in I_H},\\
 a_{ii}&=2 &&\text{for real }i,\\
 a_{ii}&\leq0 &&\text{for imaginary }i,\\
 a_{ij}&\leq0 &&\text{for }i\neq j.
\end{aligned}
\]
The evenness of the root lattice makes every imaginary diagonal entry even.

\begin{construction}[The Igusa Borcherds quiver]
Order $I_H$.  For $i<j$ place $-a_{ij}$ arrows from $i$ to $j$.  At an imaginary vertex $i$ place
\[
 g_i=1-\frac{a_{ii}}2
\]
loops.  A real vertex has loop number zero.  Denote the resulting quiver by $Q_{\Delta,H}$.
\end{construction}

\begin{proposition}[Euler form]
For dimension vectors use the Euler form
$\langle d,e\rangle=\sum_i d_ie_i-\sum_{a:i\to j}d_ie_j$.
Its symmetrization on the standard vertex basis is $A_H$:
\[
 \angles{e_i,e_j}+\angles{e_j,e_i}=a_{ij}.
\]
\end{proposition}
\begin{proof}
For $i\ne j$, the total number of arrows between the two vertices is $-a_{ij}$, so the symmetrized Euler entry is $a_{ij}$.  At $i$ the value is $2-2g_i=a_{ii}$.
\end{proof}

The inclusions $I_H\subset I_{H+1}$ identify $Q_{\Delta,H}$ with a full subquiver of $Q_{\Delta,H+1}$.  Representations supported on the old vertices form an exact extension-closed subcategory.  This gives compatible Hall embeddings.


\section{Hall realization of the positive algebra}
In this Hall construction $q$ is a prime power. Let
$\mathcal E_H=\Rep_{\mathbb F_q}(Q_{\Delta,H})$. Its twisted Hall product is
\[
 [M]*[N]
 =v^{\angles{\dim M,\dim N}}
 \sum_{[L]}
 \abs{\set{N'\subset L:N'\cong N,\ L/N'\cong M}}
 [L],
 \qquad v^2=q.
\]

The class $[0]$ is the unit: the required subobject is unique on either
side. Associativity follows by counting flags $N_1\subset N_2\subset L$
in two orders. The twist factors agree because the Euler form is
bilinear:
\[
 \langle\alpha,\beta\rangle+\langle\alpha+\beta,\gamma\rangle
 =\langle\beta,\gamma\rangle+\langle\alpha,\beta+\gamma\rangle.
\]
Thus this normalization defines the unital twisted Hall algebra on
isomorphism classes.

\begin{theorem}[Finite-window composition algebra]
\label{comp:thm:quiver-composition}
At each vertex $i$ of $Q_{\Delta,H}$, choose the one-dimensional
representation $S_i$ supported at $i$, with every loop acting by zero.
Use the generic composition algebra generated by these classes across
finite fields, with the twisted subobject-count multiplication above.
After extending coefficients from $\mathbb Z[v,v^{-1}]$ to
$\mathbb C(v)$, this is the positive quantum generalized Kac--Moody
algebra for the even Borcherds--Cartan matrix $A_H$.
The parameter in Kang--Schiffmann is $v_{\mathrm{KS}}=v^{-1}$.
Their specified semisimple perverse-sheaf construction on nilpotent
representation varieties gives the corresponding canonical basis.
\end{theorem}
\begin{proof}
This is the selected-simple composition construction of
Kang--Schiffmann, Section~2.3 and Theorem~2.2, with one selected
zero-loop simple at each vertex. Their geometric realization uses
its split graded Grothendieck group over
$\mathbb Z[v_{\mathrm{KS}},v_{\mathrm{KS}}^{-1}]$
(Theorem~4.1), with the stated inversion of parameter
\cite{KangSchiffmann,KangSchiffmannAddendum}. The selected geometric
basis objects can be semisimple, as they are for a one-loop vertex;
they are not all the individual simple perverse sheaves on the
representation stack. The theorem does not use every simple object
of the full representation category. A vertex with $c$ loops has
$q^c$ one-dimensional simple representations over $\mathbb F_q$,
whereas only the chosen zero-loop representation is the generator here.
\end{proof}

A semi-derived Hall realization of a full quantum group additionally
specifies two-periodic complexes, the invertible acyclic classes and
Cartan reduction \cite{LuBorcherdsHall,AxtellLee}.
Its coproduct, pairing and completion are separate from the positive
composition algebra used here.

\subsection{Parity and the quantized superdatum}
The Gritsenko--Nikulin denominator algebra carries a parity map on its simple-root multiset.  It is additive on the free lattice of labelled simple generators.
Projection to the Lorentzian degree can identify contributions of
different parity. Both parity dimensions must therefore be retained
in each projected root space.  Hong's construction applied to this Borcherds superdatum gives the positive part
\[
 U_{q,H}^+
\]
and the full quantized Borcherds superalgebra $U_{q,H}$.

\begin{theorem}[Ordinary and super presentations]
The ordinary quiver Hall presentation and Hong's quantum
superalgebra use the same specified Cartan matrix and labelled
simple generators. Their relation ideals and parity conventions
are different inputs. Equality of their Cartan gradings alone does
not identify their PBW characters.
\end{theorem}
\begin{proof}
The two presentations are those of the preceding Hall construction
and Hong's superalgebra construction. The difference already occurs
for a single isotropic generator. The enveloping algebra of a
one-dimensional even abelian Lie algebra is \(\mathbb C[x]\).
Declaring its generator odd only as a graded vector space leaves
a nonzero vector \(x^2\) in degree two. The enveloping algebra
of a one-dimensional odd abelian Lie superalgebra is
\(\bigwedge(\xi)\), since its defining relation is
\(2\xi^2=[\xi,\xi]=0\). Its degree-two component is zero.
Thus parity assignment to an ordinary PBW basis does not produce
the super relation ideal or its character. A comparison must
respect these relations degree by degree.
\end{proof}


\section{The real-root seed window}
The three real simple roots are
\[
 \delta_1=2f_2-f_3,
 \qquad
 \delta_2=2f_{-2}-f_3,
 \qquad
 \delta_3=f_3,
\]
with Gram matrix
\[
 A_{\mathrm{re}}=
 \begin{pmatrix}
 2&-2&-2\\
 -2&2&-2\\
 -2&-2&2
 \end{pmatrix}.
\]
Choose the order $\delta_1<\delta_2<\delta_3$.  The real-root quiver $Q_{\mathrm{re}}$ has two arrows from vertex $i$ to vertex $j$ for each $i<j$.

\begin{proposition}[Real-root Hall seed]
The symmetrized Euler form of $Q_{\mathrm{re}}$ is $A_{\mathrm{re}}$.  The generic twisted Hall composition algebra over $\mathbb C(v)$
generated by its three simple representations is the positive quantum Kac--Moody algebra with Cartan matrix $A_{\mathrm{re}}$.
\end{proposition}
\begin{proof}
The diagonal Euler entries are $2$.  The total number of arrows between distinct vertices is two, hence the symmetrized off-diagonal entries are $-2$.  The Ringel--Green theorem identifies the generic composition algebra
with the positive quantum algebra for this symmetric Cartan matrix
\cite{Ringel,Green}.
\end{proof}

The three vertices already contain the Weyl group and the reflective divisor.  Imaginary vertices add the automorphic correction.  The finite-height tower therefore begins with a completely explicit hyperbolic Hall algebra and enlarges it by one finite collection of imaginary simple roots at each stage.

\subsection{The first Ginzburg differential}
Let $a_{ij}^{(1)},a_{ij}^{(2)}:i\to j$ be the two arrows for $i<j$.  Add reverse arrows $(a_{ij}^{(r)})^*:j\to i$ of degree $-1$ and loops $t_i$ of degree $-2$.  The canonical three-Calabi--Yau completion has
\[
 d a_{ij}^{(r)}=0,
 \qquad
 d(a_{ij}^{(r)})^*=0,
\]
and
\begin{equation}\label{eq:vol3-ginzburg-differential}
 d t_i=
 \sum_{a:s(a)=i}aa^*
 -\sum_{a:t(a)=i}a^*a.
\end{equation}
The identity $d^2=0$ follows immediately.  Formula \eqref{eq:vol3-ginzburg-differential} is the derived moment-map equation for the doubled real-root quiver.

\begin{calculation}[Vertex one]
At vertex $1$,
\[
 d t_1=
 \sum_{r=1}^2
 \left(
 a_{12}^{(r)}(a_{12}^{(r)})^*
 +a_{13}^{(r)}(a_{13}^{(r)})^*
 \right).
\]
At vertex $3$,
\[
 d t_3=-\sum_{r=1}^2\left(
 (a_{13}^{(r)})^*a_{13}^{(r)}+
 (a_{23}^{(r)})^*a_{23}^{(r)}\right).
\]
The sum over vertices is $\sum_a(aa^*-a^*a)$.
It vanishes in the cyclic quotient and after matrix trace,
but generally not in the path algebra.
\end{calculation}
\begin{proof}
The outgoing arrows at vertex $1$ have targets $2$ and $3$.
The incoming arrows at vertex $3$ have sources $1$ and $2$.
The displayed formulas follow from these incidence lists.
Each arrow contributes one commutator to the sum over vertices.
The trace of a commutator vanishes, giving the asserted trace identity.
\end{proof}


\section{The three-Calabi--Yau completion}
Let $kQ_{\Delta,H}$ be the path algebra. Its algebraic
three-Calabi--Yau completion is
\[
 \Gamma_{\Delta,H}=\Pi_3(kQ_{\Delta,H}).
\]
Write $J$ for the ideal generated by all arrows in its graded quiver.
Its completed version is the separate topological dg algebra
$\widehat\Gamma_{\Delta,H}=\varprojlim_n\Gamma_{\Delta,H}/J^n$.
The algebraic Calabi--Yau theorem is used for $\Gamma_{\Delta,H}$;
a topological Calabi--Yau assertion for $\widehat\Gamma_{\Delta,H}$
requires its completed bimodule category.
This is the Ginzburg model with zero potential: a degree-minus-one
reverse arrow for every arrow and a degree-minus-two loop at each
vertex. The loop differential is the signed sum of incident
arrow--reverse-arrow products.

\begin{theorem}[Borcherds--Ginzburg projections]
For each finite quiver \(Q_{\Delta,H}\), its undeformed
three-Calabi--Yau completion \(\Gamma_{\Delta,H}\) is the
Ginzburg dg algebra with zero potential. If \(H\le H'\) and
\(Q_{\Delta,H}\) is the full subquiver on a subset of vertices,
there is a surjective dg algebra map
\[
 \pi_{H'H}:\Gamma_{\Delta,H'}\longrightarrow\Gamma_{\Delta,H}.
\]
It kills the added vertices, every arrow incident to them, the
corresponding reverse arrows, and the added degree-minus-two loops.
These maps compose strictly and extend continuously to the completed algebras.
\end{theorem}
\begin{proof}
The finite Calabi--Yau assertion is the Ginzburg model of Keller's
completion theorem \cite{KellerCY,GinzburgCY}. We concatenate
paths from left to right, so \(aa^*\) is based at the source
of \(a\). Original and reverse arrows have zero differential.
For the degree-minus-two loop \(t_i\), the differential is
\[
 d t_i=\sum_{s(a)=i}aa^*-\sum_{t(a)=i}a^*a.
\]
The projection deletes exactly the terms involving an added vertex,
leaving the same sum for the full subquiver. It therefore commutes
with the differential. It preserves path length and extends to the
completed path algebra. Successive projections kill the union of
their omitted vertices and arrows, proving strict composition.
\end{proof}

The evident inclusion of graded quivers does not give the reverse
dg map. A single old vertex has loop differential zero. Adding
an arrow \(a:i\to j\) changes that differential to \(aa^*\)
at \(i\). The map fixing the old loop therefore fails to commute
with the differential.

Write \(\Gamma_\Delta=(\Gamma_{\Delta,H},\pi_{H'H})_H\)
for this pro-system of dg algebras. It is separate from the
ordinary Hall embeddings obtained by extension by zero on
representations of a full subquiver. A direct system of
Calabi--Yau dg categories compatible with those Hall embeddings
requires additional dg functors. It does not follow from the
inclusion of the displayed Ginzburg generators.

The quantum group retains its separately specified Borel
completion. Its transition maps must preserve the Cartan generators,
coproduct, and pairing. Quotienting the positive nilpotent Lie
algebra by height and truncating a full quantum group are different
operations.

\section{Hall states and protected superstates}
The representation varieties of $Q_{\Delta,H}$ carry the
perverse-sheaf complexes used in the composition-algebra construction.
The full equivariant derived category is a larger carrier.
Already for one real vertex and dimension one, its quotient is
$B\mathbb G_m$, with $H^*(B\mathbb G_m;\mathbb Q)=\mathbb Q[c_1]$,
where $|c_1|=2$. Its total equivariant cohomology is infinite-dimensional.
Finite root multiplicities therefore refer to the specified composition
algebra, not to total equivariant cohomology of this stack.

\begin{definition}[Igusa Hall category]
For every dimension vector $\mathbf d$, let
\[
 \mathfrak R_{H,\mathbf d}=
 [\operatorname{Rep}_{\mathbf d}(Q_{\Delta,H})/G_{\mathbf d}].
\]
Use bounded equivariant constructible complexes, with coefficients
$\overline{\mathbb Q}_\ell$ away from the field characteristic.
The direct sum over dimension vectors has convolution $q_!p^*$
along the extension stack, with endpoints on the left and middle
term on the right. Flag base change proves associativity. Denote
this monoidal category by $\mathcal B_{\Delta,H}$ and set
\[
 \mathcal B_\Delta=\varinjlim_H\mathcal B_{\Delta,H}.
\]
\end{definition}

Independently, let \(\mathfrak n_{I_H}\) be the positive Lie
superalgebra generated by the selected labelled simple roots, with
their Borcherds relations. Its one-particle space is
\[
 \mathcal H_{\mathrm{1p},H}=\mathfrak n_{I_H}
       =\bigoplus_{\alpha>0}(\mathfrak n_{I_H})_\alpha.
\]
It generally has infinitely many root degrees. The bound \(H\)
restricts the simple generators; it is different from the nilpotent
quotient \(\mathfrak n_H\), which discards all root spaces above
height \(H\). The corresponding multiparticle space is the
PBW superspace of \(U(\mathfrak n_{I_H})\).

\begin{theorem}[Hall and protected-state comparison]
The composition subalgebra in the finite-window Hall theorem
is the ordinary positive quantum generalized Kac--Moody algebra.
This statement does not identify the Grothendieck group of the full
category $\mathcal B_{\Delta,H}$ with that composition algebra.  The protected one-particle superspace has signed root multiplicities
\[
 \sdim\mathcal H_{\mathrm{1p},H,\alpha}
 =\operatorname{sdim}(\mathfrak n_{I_H})_\alpha.
\]
Taking the union of labelled windows and then completing by total
root height gives the PBW character
\[
 \sch\Fock(\nIg)
 =\prod_{\alpha>0}(1-e^{-\alpha})^{-\sdim\gIg_\alpha},
\]
and its Weyl-shifted inverse is $\Dmon(2Z)$.  The composition canonical basis and the protected root presentation use the specified simple-root indexing and Cartan grading; the parity refinement belongs to the Borcherds superdatum.
\end{theorem}
\begin{proof}
The composition-algebra assertion is the finite-window theorem. The protected assertion is PBW for the Borcherds Lie superalgebra.  Proper height passes both constructions to the limit.  The Gritsenko--Nikulin denominator theorem gives the final product identity.
\end{proof}

\section{Extension determinants and critical orientations}
Let \(\mathcal C\) be a three-Calabi--Yau dg category over
\(\mathbb C\), with its chosen trace. For families of objects
\(A,B\) over a base \(T\), let \(C(A,B)\) denote the relative
derived Hom complex on \(T\). Suppose these complexes are perfect,
and that the trace gives base-change-compatible duality
\[
 C(B,A)\simeq C(A,B)^\vee[-3].
\]
Put \(L(A,B)=\det C(A,B)\) and \(D(A)=L(A,A)\).
The determinant convention is
\(\det C=\bigotimes_i(\det C^i)^{(-1)^i}\), with symmetry sign
specified by \(\chi(C)\bmod2\).

\begin{theorem}[Extension determinant]
For an exact triangle \(A\to F\to B\to A[1]\), there is a
base-change-compatible isomorphism
\[
 D(F)\simeq D(A)\otimes D(B)\otimes L(A,B)^{\otimes2}.
\]
The determinant comparisons obtained by refining a flag agree, with
the stated determinant symmetry signs.
\end{theorem}
\begin{proof}
Apply derived Hom to the triangle in each variable. Determinant
additivity gives
\[
 D(F)\simeq L(B,A)\otimes D(A)\otimes D(B)\otimes L(A,B).
\]
Dualization and the odd shift each invert the determinant line.
Thus Calabi--Yau duality identifies \(L(B,A)\) with \(L(A,B)\).
Reorder by determinant symmetry to obtain the formula. For a full
flag, all iterated decompositions identify with the same ordered
tensor product of \(C(A_i,A_j)\), for its successive quotients
\(A_i\). Determinant additivity for refined filtrations proves
coherence. Each operation commutes with base change.
\end{proof}

A square-root orientation consists of lines \(o_A\) and
isomorphisms \(o_A^{\otimes2}\simeq D(A)\). Extension
compatibility requires maps
\[
 o_A\otimes o_B\otimes L(A,B)\longrightarrow o_F
\]
whose squares are the determinant comparisons. Choosing coordinates
on a smooth critical chart trivializes its ambient canonical line.
Hall convolution requires transition maps compatible with the cross
line and the zero object. On three-step flags, the two orientation
maps can differ by an element of \(\mu_2\), even when their
squares agree. That sign must be trivialized coherently.

\begin{example}[A nontrivial cross line]
Let \(X\) be a smooth projective Calabi--Yau threefold, with a
chosen volume form, and let \(x\in X\). On
\(X\times T\), where \(T=\mathbb P^1\), take
\[
 A=\mathcal O_X\boxtimes\mathcal O_T(1),\qquad
 B=\mathcal O_x\boxtimes\mathcal O_T.
\]
Then \(C(A,B)=\mathcal O_T(-1)[0]\), and relative Serre
duality gives \(C(B,A)=\mathcal O_T(1)[-3]\).
Thus \(L(A,B)=L(B,A)=\mathcal O_T(-1)\).
The self-Hom complexes of \(A\) and \(B\) have constant
determinant lines, since the factors \(\mathcal O_T(1)\)
cancel in the first self-Hom. Nevertheless
\[
 D(A\oplus B)=D(A)\otimes D(B)\otimes\mathcal O_T(-2).
\]
The orientation cross factor is \(\mathcal O_T(-1)\).
A sign change cannot replace this line of nonzero degree.
\end{example}

\begin{corollary}[Oriented critical coefficients]
Suppose the chosen critical charts have square-root orientation
maps on overlaps and extension flags, with unit and flag coherences.
Suppose their vanishing-cycle comparison maps preserve these data.
Then the local vanishing-cycle complexes descend to the specified
oriented coefficient object. Their Hall convolution is associative
whenever its pullback and pushforward maps are defined and satisfy
base change on the flag diagrams.
\end{corollary}
\begin{proof}
The overlap cocycle identifies the local coefficient complexes by
descent. On a three-step flag, orientation coherence identifies
the two Thom--Sebastiani coefficient maps. Proper or compact-support
base change, on the stipulated domain, identifies the two composites
of convolution. The zero flag gives the unit equation.
\end{proof}

For the quiver charts, the remaining construction is an equivariant
system of square roots and coefficient maps on the actual
representation and extension stacks, compatible with changes of
height. For compact \(S\times E\), the reduced complexes and
the common elliptic torsor must also be retained. A coordinate
volume on an affine chart specifies neither descent system by itself.


\chapter{The quantum root datum}
\label{comp:ch:quantum-roots}

The labelled presentation also admits a quantum Hopf algebra.
Its Cartan lattice, parity signs and Borel completion are part of
the construction. The canonical pairing produces a root tensor;
a separate Cartan weight operator turns it into the intertwiner
on the specified representation domain.

\section{The integral colored datum}
Retain the labelled simple-root set \(I_\Delta\), including its
multiplicities and parity \(p(i)\). For a finite sublevel
\(I_H\), use the free lattice
\[
 Q_H=\bigoplus_{i\in I_H}\mathbb Z\epsilon_i,
 \qquad \pi(\epsilon_i)=\alpha_i\in L_{II},
 \qquad a_{ij}=(\alpha_i,\alpha_j).
\]
The sign bicharacter is
\(\theta(\epsilon_i,\epsilon_j)=(-1)^{p(i)p(j)}\).
It lives on the labelled lattice. Since \(\pi\) can identify
opposite parities, it need not descend to \(L_{II}\).

The matrix \(A_H\) is integral and symmetric, with real diagonal
two, imaginary diagonal nonpositive and even, and off-diagonal
entries nonpositive. Odd simple generators are imaginary. These
are the hypotheses for the colored quantum Borcherds construction
\cite[Section~2]{Hong98}. Work over \(\mathbb C(q)\), where
\(q\) is indeterminate, so no root-of-unity specialization is implied.

For the extended Cartan, take independent \(h_i,d_i\) with
\[
 (h_i,h_j)=a_{ij},\qquad(h_i,d_j)=\delta_{ij},
 \qquad(d_i,d_j)=0.
\]
Its form is nondegenerate. The simple root as a Cartan weight is
\(x\mapsto(h_i,x)\); the \(d_i\) separate repeated labelled
roots. The derived quantum algebra uses only the torus generated
by the \(h_i\), together with the positive and negative generators.

\section{The reduced Cartan and coproduct}
The derived algebra is generated by \(E_i,F_i,K_h\), with
\(h\in\bigoplus_i\mathbb Zh_i\), \(|E_i|=|F_i|=p(i)\), and
\[
 K_hK_{h'}=K_{h+h'},\qquad
 K_hE_iK_h^{-1}=q^{(h,h_i)}E_i,\qquad
 K_hF_iK_h^{-1}=q^{-(h,h_i)}F_i.
\]
Its mixed relation is
\[
 E_iF_j-(-1)^{p(i)p(j)}F_jE_i
 =\delta_{ij}\frac{K_{h_i}-K_{h_i}^{-1}}{q-q^{-1}}.
\]
For real \(i\ne j\), the positive Serre relation is
\[
 \sum_{r=0}^{1-a_{ij}}(-1)^r
 {1-a_{ij}\brack r}_q E_i^{1-a_{ij}-r}E_jE_i^r=0,
\]
with the corresponding relation for \(F\). If \(a_{ij}=0\),
the two generators supercommute. Here
\([n]_q=(q^n-q^{-n})/(q-q^{-1})\), and the bracket denotes
the associated Gaussian binomial coefficient. All tensor products
of superalgebras use the Koszul multiplication sign.

The coproduct and counit are
\[
\begin{aligned}
 \Delta E_i&=E_i\otimes1+K_{h_i}\otimes E_i,
 &\epsilon(E_i)&=0,\\
 \Delta F_i&=F_i\otimes K_{h_i}^{-1}+1\otimes F_i,
 &\epsilon(F_i)&=0,\\
 \Delta K_h&=K_h\otimes K_h,
 &\epsilon(K_h)&=1.
\end{aligned}
\]
The antipode is \(S(E_i)=-K_{h_i}^{-1}E_i\),
\(S(F_i)=-F_iK_{h_i}\), and \(S(K_h)=K_{-h}\).

\begin{theorem}[Quantization with the Lorentzian Cartan]
Impose \(K_h=1\) for
\[
 h\in\ker\!\left(\bigoplus_{i\in I_H}\mathbb Zh_i
                 \longrightarrow L_{II},\ h_i\longmapsto\alpha_i\right).
\]
This is a central Hopf quotient. Once \(I_H\) contains the three
real generators, its torus is \(\mathbb C(q)[L_{II}]\), and it has
triangular decomposition
\[
 U_{q,H}^{\mathrm{red}}
 \cong U_{q,H}^-\otimes\mathbb C(q)[L_{II}]\otimes U_{q,H}^+.
\]
The positive and negative labelled PBW dimensions and parities agree
with those of the classical Borcherds presentation. The inclusions
of finite labelled sets give a quantum algebra
\(U_q^{\mathrm{red}}(\gIg)\) with these properties.
\end{theorem}
\begin{proof}
Hong's presentation and triangular decomposition apply to the
extended datum \cite[Propositions~2.2--2.3 and~3.1]{Hong98}.
A Cartan vector in the displayed kernel pairs to zero with every
\(h_j\). Its torus element is therefore central. Moreover
\[
 \Delta(K_h-1)=(K_h-1)\otimes K_h+1\otimes(K_h-1),
\]
so these relations generate a Hopf ideal. In the triangular
decomposition the quotient acts only on the middle group algebra.
Its quotient is the group algebra of the image lattice, which is
\(L_{II}\) because the real roots form a basis.

The labelled positive and negative relations remain unchanged.
Hong's character theorem identifies their graded dimensions with
the classical ones. Relations involving an added generator have
positive degree in that generator, so they impose no new relation
on an older positive or negative subalgebra. Triangular decomposition
then gives the injective transition maps. In a fixed positive
Lorentzian degree, proper height permits only finitely many labelled
words. The finite character identity therefore survives both the
union and the projection to that degree, in each parity.
\end{proof}

The full quantum algebra has not been completed by taking a product
of all positive and negative root spaces. The discontinuity in
Proposition~\ref{comp:prop:two-sided-product} persists as a warning
against that topology. The Borel completion used for the canonical
tensor has a different definition.

\section{The canonical tensor and the weight operator}
For a finite labelled set, let \(U\) be the extended quantum
algebra. If \(\beta\in Q_H^+\), let
\[
 U^{+,\beta}=\bigoplus_{\gamma\nleq\beta}U^+_\gamma,
 \qquad \widehat U=\varprojlim_\beta U/U\,U^{+,\beta}.
\]
The order is coordinatewise on the labelled positive lattice.
These are left-ideal neighbourhoods, not two-sided root-height
quotients. Hong constructs the algebra structure and its tensor
completion in this topology \cite[(7.2)]{Hong98}.

The nondegenerate Borel pairing gives a canonical element
\(C_\beta\in U^+_\beta\otimes U^-_{-\beta}\).
Put \(h_\beta=\sum_i\beta_i h_i\) and \(K_\beta=K_{h_\beta}\).
The completed element is
\[
 C=\sum_{\beta\in Q_H^+}\theta(\beta,\beta)
       q^{(h_\beta,h_\beta)}
       (K_\beta^{-1}\otimes K_\beta)C_\beta.
\]
It satisfies a twisted intertwining identity. The twist \(\Phi\)
fixes pairs of torus elements and has
\[
\begin{aligned}
 \Phi(E_i\otimes1)&=E_i\otimes K_i,&
 \Phi(1\otimes E_i)&=K_i\otimes E_i,\\
 \Phi(F_i\otimes1)&=F_i\otimes K_i^{-1},&
 \Phi(1\otimes F_i)&=K_i^{-1}\otimes F_i.
\end{aligned}
\]
Here \(K_i=K_{h_i}\). On weight modules for which the weight
pairing is defined, set
\[
 Z(v_\lambda\otimes w_\mu)=q^{(\lambda,\mu)}
                              v_\lambda\otimes w_\mu.
\]
Adjoin the required powers of \(q\) if the chosen weight lattice
has rational pairing. Restrict to tensor products on which the
series for \(C\) and its inverse are defined. For example, a
highest-weight bound on the first factor makes each vector receive
only finitely many positive-root terms.

\begin{theorem}[The quantum \(R\) operator]
On this weight-module domain,
\[
 R=Z^{-1}C
\]
is invertible and satisfies \(R\Delta(a)=\Delta^{\mathrm{op}}(a)R\)
and the Yang--Baxter equation. The opposite coproduct uses the
super flip. These conclusions hold for the full labelled system
on modules satisfying the same finiteness conditions.
\end{theorem}
\begin{proof}
Nondegeneracy of the Borel pairing is
\cite[Theorem~4.6]{Hong98}. The canonical tensor satisfies
\(C\Delta(a)=\Phi(\Delta^{\mathrm{op}}a)C\) and the two
compatible coproduct identities of \cite[Propositions~7.4--7.5 and Theorem~7.6]{Hong98}.
Conjugation by \(Z\) realizes \(\Phi\). Multiplying by
\(Z^{-1}\) gives the ordinary intertwining identity. The two
coproduct identities give the Yang--Baxter equation, as in
\cite[Theorem~7.7]{Hong98}. The finiteness hypothesis permits these
identities to act on the stated tensor products. In the full system,
every contributing finite set of labelled degrees lies in a finite
window, where the same identities hold.
\end{proof}

The Cartan factor is necessary even for one real root. On the
standard two-dimensional \(U_q(\mathfrak{sl}_2)\)-module,
choose $v_+,v_-$ with $Kv_\pm=q^{\pm1}v_\pm$, $Ev_-=v_+$,
$Fv_+=v_-$, and $Ev_+=Fv_-=0$. Then
\[
 C=1-(q-q^{-1})E\otimes F,\qquad
 Z=\operatorname{diag}(q^{1/2},q^{-1/2},q^{-1/2},q^{1/2}).
\]
With the coproduct above, \(Z^{-1}C\) intertwines it with its
opposite; \(C\) alone does not. This example also fixes the sign
of the root term. The reduced Cartan quotient has the Hopf algebra
constructed above. Its braided representation theory must retain
weight pairings and a domain on which the same operator descends.

A geometric Hall comparison requires a coproduct, Hopf pairing and
module functor compatible with this quantization. A current or
fusion construction requires a further comparison of its supported
operations with the stated weight-module braiding. Neither follows
from equality of denominator characters.

\chapter{Fredholm determinants and compact comparisons}
\label{comp:ch:root-operators}

A formal determinant becomes an analytic operator determinant only
when both parity dimensions are summable. The minimal simple-root
presentation gives such a bound through an absolute theta estimate
and the free tensor algebra. This constructs the root Fredholm
model before asking for its compact geometric realization.

\section{Finite determinants}
For a positive height bound \(H\), set
\[
 \mathcal V_H=\bigoplus_{h(\alpha)\leq H}\gIg_\alpha,
 \qquad \mathscr D_H|_{\gIg_\alpha}=(1-e^{-\alpha})\operatorname{id}.
\]
All root spaces here are positive. Their two parity dimensions are
those of the fixed minimal presentation.

\begin{proposition}[Finite Berezinian and formal limit]
In the total-height completion,
\[
 \operatorname{Ber}(\mathscr D_H)
   =\prod_{h(\alpha)\leq H}(1-e^{-\alpha})^{\operatorname{sdim}\gIg_\alpha},
 \qquad
 e^{-\rho}\lim_H\operatorname{Ber}(\mathscr D_H)=\Dmon(2Z).
\]
Pullback by \(\jmath(Z)=Z/2\) sends the last expression to
\(\Dmon(Z)\), whose square is \(\ChiTen(Z)\).
\end{proposition}
\begin{proof}
On a finite super vector space, the Berezinian is the even determinant
divided by the odd determinant. This gives the product. Each formal
coefficient uses only finitely many positive-height roots, so it
stabilizes as \(H\) increases. The denominator theorem supplies
the Weyl prefactor and the coordinate substitution.
\end{proof}

\section{An absolute bound for both root parities}
Let \(M(t)\) count the labelled simple generators by their ordinary
dimension, with even and odd generators both counted positively.
For the minimal presentation,
\[
 M(t)=3t+\frac{27t^2}{1-t^2}
       +\sum_{\substack{a\in\mathcal C\cap L_{II}\\(a,a)<0}}
                                      |m(a)|t^{h(a)}.
\]
The second term counts nine even generators at every multiple of
each of the three height-two null rays.

\begin{lemma}[Theta majorant]
The series \(M(t)\) converges for \(0<t<1\), and
\(M(t)\to0\) as \(t\to0^+\).
\end{lemma}
\begin{proof}
Specialize the root-torus variables to
\[
 Q=S=t^2,\qquad R=t^{-1}.
\]
A monomial of root degree \(\alpha(n,\ell,m)\) has exponent
\(2n-\ell+2m=h(\alpha)\). A term of a genus-two theta
constant, indexed by \(x,y\in\mathbb Z+\{0,1/2\}\), has
absolute value
\[
 t^{x^2+y^2-xy},\qquad
 x^2+y^2-xy=\tfrac12((x-y)^2+x^2+y^2)
             \geq\tfrac12(x^2+y^2).
\]
Each absolute theta series is therefore Gaussian and converges.
The product of the ten absolute series bounds the sum of absolute
Fourier coefficients of \(\Dtheta\), including possible
cancellations in its actual coefficients. Since
\(m(a)=-c_{\rho+a}\), the last sum in \(M(t)\) is bounded
by this product divided by \(64t^{h(\rho)}\), where
\(h(\rho)=3/2\). This proves convergence at every fixed
\(0<t<1\). All terms of \(M\) have positive integral height.
Dominated convergence using any fixed \(t_*<1\) then proves
\(M(t)\to0\).
\end{proof}

\begin{theorem}[Full root-dimension bound]
There exist \(C_0,C_1>0\) such that
\begin{equation}\label{eq:vol3-root-growth}
 \dim\gIg_\alpha\leq C_0\exp(C_1h(\alpha))
\end{equation}
for every positive root. More strongly, for any \(t_0\in(0,1)\)
with \(M(t_0)<1\),
\[
 \sum_{\alpha>0}\dim\gIg_\alpha\,t_0^{h(\alpha)}
 \leq\frac{1}{1-M(t_0)}-1.
\]
Both parity dimensions are included in this estimate.
\end{theorem}
\begin{proof}
Let \(V\) be the superspace of labelled simple generators. The
tensor algebra has ordinary dimension series
\[
 \sum_{r\geq0}\operatorname{Hilb}(V^{\otimes r};t_0)
       =\sum_{r\geq0}M(t_0)^r=\frac1{1-M(t_0)}.
\]
It surjects onto \(U(\nIg)\). PBW embeds \(\nIg\) in that
enveloping algebra. All maps preserve height and parity, so dimensions
can only decrease. Removing the unit gives the stated sum. Each
summand is bounded by its right side, giving the exponential estimate
with \(C_1=-\log t_0\). The lemma ensures that such \(t_0\)
exists.
\end{proof}

A signed multiplicity bound alone would not prove this theorem.
At the first timelike root it bounds 64, whereas the full root
space has dimension 122. The tensor-algebra argument starts from
the specified simple generators and controls all their brackets.

\section{The two root Hilbert spaces}
Choose a homogeneous orthonormal basis in every root space. Define
\[
 \mathcal H_\Delta=\widehat{\bigoplus}_{\alpha>0}\gIg_\alpha.
\]
There is also a minimal virtual space
\[
 V^{\mathrm{vir}}_\alpha
  =\mathbb C^{\max(\operatorname{sdim}\gIg_\alpha,0)
                  \mid\max(-\operatorname{sdim}\gIg_\alpha,0)},
 \qquad
 \mathcal H^{\mathrm{vir}}_\Delta
       =\widehat{\bigoplus}_{\alpha>0}V^{\mathrm{vir}}_\alpha.
\]
The hats here mean Hilbert direct sums. The first carrier retains
both root parities; the second retains only their difference and
has no supplied root bracket.

On the denominator period domain, put \(Z_{\mathrm{root}}=2Z\).
Thus \(Q=e^{4\pi i\tau}\), \(R=e^{4\pi iz}\),
\(S=e^{4\pi i\sigma}\), and use the additive exponent
\[
 \alpha(Z)=-4\pi i(n\tau+\ell z+m\sigma),\qquad
 e^{-\alpha(Z)}=Q^nR^\ell S^m.
\]
Let \(\mathfrak C_{\mathrm{abs}}\) be the domain where
\(\Re\delta_i(Z)>0\) for all three real simple roots and
\[
 \sum_{\alpha>0}\dim\gIg_\alpha\,
                  |e^{-\alpha(Z)}|<\infty.
\]

\begin{theorem}[Nonempty trace-class domain]
For every direction \(Z_0\) with
\(c=\min_i\Re\delta_i(Z_0)>0\), the ray \(TZ_0\) lies
in \(\mathfrak C_{\mathrm{abs}}\) for all sufficiently large
\(T>0\). On both Hilbert superspaces the operator
\[
 K_Z|_\alpha=e^{-\alpha(Z)}\operatorname{id}
\]
is trace class in each parity. The family is holomorphic in trace
norm on the interior of its convergence domain.
\end{theorem}
\begin{proof}
For \(\alpha=\sum_i b_i\delta_i\), with \(b_i\geq0\),
\(\Re\alpha(TZ_0)\geq Tc\,h(\alpha)\). Choose \(T\)
so that \(e^{-Tc}\leq t_0\) in the preceding theorem. Its
summable bound dominates the trace norm of \(K_Z\). The virtual
space has no larger dimension in any degree, so the same estimate
applies there. At any interior point of the convergence domain,
choose a nearby convergent point with all three real simple
exponents smaller. Its absolute series dominates a neighbourhood.
The diagonal holomorphic entries then give a trace-norm holomorphic
family.
\end{proof}

\section{Skew doubling and the Fredholm product}
Put \(D_Z=1-K_Z\). The real parts are positive on the chosen
domain, so no eigenvalue of \(K_Z\) is one. Its even and odd
Fredholm determinants are therefore nonzero. On either root
Hilbert superspace form, separately in each parity,
\[
 \mathbb A_{Z,\bar p}=
 \begin{pmatrix}0&D_{Z,\bar p}\\-D_{Z,\bar p}^{\mathsf t}&0\end{pmatrix}.
\]
The transpose uses the chosen basis and its dual. Order the finite
root truncations by height and then by a fixed order within each
degree. Normalize each Pfaffian by its value at \(K=0\).

\begin{theorem}[Oriented Fredholm realization]
The normalized finite Pfaffians converge, and their even-to-odd
ratio satisfies
\[
 \operatorname{sPf}_F(\mathbb A_Z)
 =\frac{\det_F(D_Z|_{\bar0})}{\det_F(D_Z|_{\bar1})}
 =\prod_{\alpha>0}(1-e^{-\alpha(Z)})^{\operatorname{sdim}\gIg_\alpha}.
\]
For both the full and the minimal virtual root carrier,
\[
 e^{-\rho}\operatorname{sPf}_F(\mathbb A_Z)=\Dmon(2Z)
\]
in the Gritsenko--Nikulin period convention.
\end{theorem}
\begin{proof}
For a finite square matrix \(D\), the Pfaffian of
\(\left(\begin{smallmatrix}0&D\\-D^t&0\end{smallmatrix}\right)\)
is \((-1)^{N(N-1)/2}\det D\). Normalizing at \(D=1\)
removes that sign. Hence the paritywise normalized Pfaffians equal
their determinants at every finite height. The truncations of
\(K_Z\) converge in trace norm, so determinant continuity gives
the Fredholm limit. The odd determinant is nonzero on the domain,
allowing the ratio. Both carriers have the same superdimension in
every degree, so they give the same product. The denominator
identity supplies its Weyl shift.
\end{proof}

This is a skew operator model with a specified polarization and
relative Pfaffian normalization. It does not supply a first-order
differential operator or its principal symbol. A compact geometric
family requires its own domain, symbol, line comparison and
orientation transport, together with an intertwiner to the root model.

\section{What the automorphic character determines}
\begin{proposition}[Information retained by the product]
With the cusp and automorphy law fixed, the normalized product
records its signed exponents, Weyl prefactor, divisor and modular
multiplier. It does not determine a Lie bracket, the two separate
parity dimensions, extension correspondences, or an orientation
comparison with a geometric determinant line.
\end{proposition}
\begin{proof}
The recorded data are the exponents, prefactor and transformation
law. To see the lost information concretely, put even vectors
\(x,y\) in degree one and \(z\) in degree two. The abelian
bracket and the Heisenberg bracket \([x,y]=z\) give nonisomorphic
Lie algebras with the same PBW character
\[
 \frac1{(1-t)^2(1-t^2)}.
\]
An added even--odd pair leaves every signed exponent unchanged.
Neither character calculation specifies an extension correspondence.
For the line ambiguity, the two square-root pairs on
\(\mathbb C^\times\), with square isomorphisms given by \(1\)
and by the coordinate \(t\), have the same square line. They
are not isomorphic as root pairs, since no unit of
\(\mathbb C[t,t^{-1}]\) has square \(t\). A comparison must
retain these structures separately.
\end{proof}

\begin{recognition}[Primitive recognition from characters]
Let \(\mathfrak p\) satisfy the proper-height finiteness hypotheses
of Theorem~\ref{igusa:primitive-presentation-map}, and let
\(\Psi:\nIg\twoheadrightarrow\mathfrak p\) be a graded,
parity-preserving Lie map. It is an isomorphism if the two
parity-refined PBW series agree. More generally, for a surjection
between positive Lie superalgebras whose root spaces have
prescribed pure parity, equality of their signed PBW characters
also suffices. Their induced height completions, Chevalley complexes
and current envelopes are then isomorphic.
\end{recognition}
\begin{proof}
The first assertion is the universal presentation criterion.
For the second, the logarithm of the PBW product determines signed
primitive dimensions inductively in height: contributions from
proper multiples use smaller degrees and can be subtracted.
Pure parity recovers ordinary dimensions from these signed
numbers. Degreewise surjectivity and finite dimension give
injectivity. All the other maps are induced functorially.
\end{proof}

The pure-parity criterion does not apply to the first timelike
Igusa root, where both parities occur. Its comparison uses both
parity dimensions.

\subsection{The automorphic line and its reflection action}
The scalar product, with its Weyl prefactor, continues to a section
of \(\mathcal L^5\otimes\nu_5\), where \(\mathcal L\) is
the weight-one modular line. This constructs its automorphic
transition functions. Operator transition maps are additional:
scalar determinant equality does not identify Hilbert spaces,
operator domains or their skew families.

Proposition~\ref{comp:prop:reflection-signs} specifies the linear
and weight-five slash actions separately.


An orientation comparison with a geometric skew family must specify
which of these actions it intertwines, together with the weight-five
line. The monodromy of a root line is not determined solely by the
vanishing order of one scalar section.

\section{Reduced curve counting on \texorpdfstring{$K3\times E$}{K3 x E}}
Let \(S\) be a smooth projective K3 surface and \(E\) an
elliptic curve. For a nonzero effective primitive class
\(\beta_h\in H_2(S,\mathbb Z)\), with \(\beta_h^2=2h-2\),
choose \(\beta_h^\vee\in H^2(S,\mathbb Q)\) satisfying
\(\int_{\beta_h}\beta_h^\vee=1\). Define
\[
 N_{g,h,d}=\int_{[\overline M_{g,1}^{\bullet}
       (S\times E,(\beta_h,d))]^{\mathrm{red}}}
       \operatorname{ev}_1^*(\beta_h^\vee\otimes[\mathrm{pt}]_E).
\]
The bullet permits disconnected domains with no collapsed connected
components; \(2g-2\) denotes minus their Euler characteristic.
The insertion fixes the reduced incidence convention. Reduced
deformation invariance makes the result depend only on $g,h,d$.
As $h$ varies, use any pairs $(S,\beta_h)$ with the stated square. Put
\[
 \mathcal Z(u,q,s)=\sum_{g,h,d\geq0}
                 N_{g,h,d}u^{2g-2}q^{h-1}s^{d-1}.
\]
\begin{theorem}[Primitive Igusa cusp-form theorem]
The primitive reduced series is the Laurent expansion in $u$ of
\[
 \mathcal Z(u,q,s)=
   \left.-\ChiTen(q,r,s)^{-1}\right|_{r=e^{iu}}.
\]
For fixed $q,s$ degrees, this is the expansion of the corresponding
rational function of $r$ at $r=1$.
Equivalently, in theta normalization,
\[
 \mathcal Z=-2^{12}\Dtheta^{-2}.
\]
\end{theorem}
\begin{proof}
Oberdieck and Pixton identify the primitive reduced partition function with the Laurent expansion of $-\ChiTen^{-1}$ \cite{OP}; their normalization is the monic Borcherds product $\ChiTen$. The second formula is the identity $\ChiTen=2^{-12}\Dtheta^2$.
\end{proof}

The theorem identifies the numerical generating series in its stated
primitive sector. The Borcherds root algebra supplies a separate
algebraic character to compare with that series.

\begin{corollary}[Exact character comparison]
Let $\jmath:\HHtwo\to\HHtwo$ be the half-coordinate map $\jmath(Z)=Z/2$, and let
\[
 \Fock(\nIg)=
 \Sym(\nIg_{\bar0})\tensor\bigwedge(\nIg_{\bar1}).
\]
Then
\[
 \mathcal Z^{\red}_{S\times E}
 =-\jmath^*\!\left(
 e^{2\rho}\sch\bigl(\Fock(\nIg)^{\tensor2}\bigr)
 \right)
 =-\ChiTen^{-1}.
\]
\end{corollary}
\begin{proof}
Put
\[
 P=\prod_{\alpha>0}(1-e^{-\alpha})^{\sdim\gIg_\alpha}.
\]
The denominator theorem gives $\den(\gIg)=e^{-\rho}P=\Dmon(2Z)$, while PBW gives $\sch\Fock(\nIg)=P^{-1}$.  Hence
\[
 e^{2\rho}\sch\bigl(\Fock(\nIg)^{\tensor2}\bigr)
 =\den(\gIg)^{-2}.
\]
Pullback by $\jmath$ sends $\den(\gIg)$ to $\Dmon(Z)$, whose square is $\ChiTen$.  The reduced incidence convention supplies the overall minus sign in the Oberdieck--Pixton theorem.
\end{proof}

The equality compares the completed root-Fock supercharacter with the
numerical reduced generating series. Its proof uses the denominator
identity and the half-coordinate pullback. A map from compact geometric
states to this Fock space, preserving operators and their trace,
is an additional comparison.


\section{The eight diagonal-divisor forms}
Cl\'ery and Gritsenko classify eight genus-two diagonal-divisor
forms \cite{GC}. The table records their groups, weights and
character orders. Half-integral weights are interpreted on the
corresponding multiplier covers. These eight rows are not a list
of cyclic CHL orders. The reflective presentations in
\cite{GN97,GN98a,GN98b,GovindarajanKrishna2009,Govindarajan2011,GovindarajanSamanta2019}
have their own specified Jacobi inputs, lattices and chambers.
\[
\begin{array}{c|c|c|c}
\text{form}&\text{group}&\text{weight}&\text{character order}\\ \hline
\Delta_5&\operatorname{Sp}_4(\Z)&5&2\\
\Delta_2&\Gamma_2&2&4\\
\Delta_1&\Gamma_3&1&6\\
\Delta_{1/2}&\Gamma_4&1/2&8\\
\nabla_3&\Gamma^{(2)}_0(2)&3&2\\
\nabla_2&\Gamma^{(2)}_0(3)&2&2\\
\nabla_{3/2}&\Gamma^{(2)}_0(4)&3/2&4\\
Q_1&\Gamma_2(2)&1&4
\end{array}
\]

\begin{proposition}[Product models and presentation refinements]
Fix one row with an integral-exponent product
\[
 F(Z)=e^{-\rho_F(Z)}\prod_{\gamma>0}(1-e^{-\gamma(Z)})^{c_\gamma}
\]
in its declared cusp coordinates and proper positive cone. Its
minimal virtual root superspace has superdimension \(c_\gamma\)
and gives this formal determinant after the Weyl prefactor. On any
domain where \(\sum_\gamma|c_\gamma e^{-\gamma(Z)}|<\infty\)
and the factors are nonzero, it also gives a Fredholm Berezinian.
If a symmetrizable Borcherds superpresentation for the same product
is supplied, the quiver, Calabi--Yau and quantum constructions above
apply to that datum at their respective stated domains.
\end{proposition}
\begin{proof}
Choose the minimal even or odd vector space according to the sign
of \(c_\gamma\). Proper height proves formal local finiteness;
the summability assumption is exactly the trace-norm condition for
its diagonal operator. Their determinants give the displayed
product. A supplied presentation specifies the labelled Cartan
matrix and parity, which are the inputs to the later constructions.
The ordinary quiver Hall algebra uses the even relation ideal;
the quantum superalgebra uses the supplied parity relations. The
shared matrix alone does not identify those two algebras.
\end{proof}

The Igusa row has the full root-space summability proof given above.
For another row, an absolute coefficient estimate must be checked
against its actual product or simple presentation. Neither a
classification by divisor nor the equality of one scalar product
constructs its root bracket, quantum pairing or compact Hall map.

\section{The collection of comparison targets}
\begin{definition}[Igusa comparison targets]
Associated with the fixed minimal presentation are
\[
 \mathfrak K_\Delta=
 \left(\Gamma_\Delta,\mathcal B_\Delta,\gIg,
       (V_H,\mathcal A_H,\mathcal F_H)_H,
       U_q^{\mathrm{red}}(\gIg),\mathbb A_Z\right).
\]
These are respectively the pro-system of finite Calabi--Yau
completions, the ordinary quiver Hall categories, the root Lie
superalgebra, the current and Ran constructions, the quantum Hopf
superalgebra, and the oriented skew family on the root Hilbert space.
The finite BV cotangent models and the classical local BF theories
provide separate field-theoretic comparison targets.
\end{definition}

\begin{theorem}[Scalar invariants of the root constructions]
Put
\[
 P=\prod_{\alpha>0}(1-e^{-\alpha})^{\operatorname{sdim}\gIg_\alpha}.
\]
The ordinary Chevalley coalgebra has Euler supercharacter \(P\),
the positive PBW/Fock superspace has supercharacter \(P^{-1}\),
and the skew Fredholm model has super-Pfaffian \(P\) on its
convergence domain. Consequently
\[
 e^{-\rho}P=\den(\gIg)=\Dmon(2Z),\qquad
 e^{\rho}\operatorname{sch}\operatorname{Fock}(\nIg)
     =\den(\gIg)^{-1},
\]
and in the primitive reduced incidence convention,
\[
 \mathcal Z^{\mathrm{red}}_{S\times E}
    =-\bigl(\jmath^*\den(\gIg)\bigr)^{-2}
    =-\ChiTen^{-1}.
\]
\end{theorem}
\begin{proof}
The ordinary bar--Chevalley comparison gives the first character;
PBW gives its reciprocal. The Fredholm theorem gives the same
signed product on its trace-class domain. The denominator theorem
and the primitive reduced curve-counting theorem then give the
remaining identities with their stated Weyl and coordinate factors.
\end{proof}

The constructions retain distinct operations. The minimal root
datum supplies a bracket and both parities. The quiver supplies
extension correspondences and a separate ordinary composition
algebra. The current envelope supplies local fields, and its
associative Ran envelope supplies a bar in a specified tensor
category. Hong's Borel pairing supplies an \(R\) operator on its
weight-module domain. The skew doubling supplies a Pfaffian ratio,
while the automorphic form supplies a global modular section.
These scalar identities preserve the common denominator without
identifying these carriers.

For a compact realization, the first maps are the reduced critical
coefficient maps on common-torsor extensions and full flags. They
must preserve the additive physical charge and its polarized Gram
refinement. Primitive representatives must satisfy the simple-root
relations and generate with the required parity dimensions, so that
Theorem~\ref{igusa:primitive-presentation-map} applies. A comparison
of doubles additionally preserves a coproduct and full Hopf pairing.
A local--wrapped comparison preserves the supported mixed operations.
Finally an operator and state comparison must identify the relevant
orientation lines, sections, domains and sewing maps. The classical
cotangent action, the finite quantum master equation and the reduced
scalar formula leave these geometric constructions as separate
mathematical problems.
